\documentclass[11pt,a4paper]{article}

\usepackage[T1]{fontenc}
\usepackage[utf8]{inputenc}
\usepackage[ngerman]{babel}
\usepackage{amsmath,amssymb,amsthm}
\usepackage[a4paper,margin=2.5cm]{geometry}

\title{Restklassenmittel nach Schalenreduktion und die numerische Hypothese \(B>C>A>E\)}
\author{Kollaborative Ausarbeitung}
\date{\today}

\newtheorem{satz}{Satz}
\newtheorem{lemma}[satz]{Lemma}
\newtheorem{korollar}[satz]{Korollar}
\theoremstyle{definition}
\newtheorem{definition}[satz]{Definition}
\newtheorem{hypothese}[satz]{Hypothese}
\newtheorem{beispiel}[satz]{Beispiel}
\theoremstyle{remark}
\newtheorem{bemerkung}[satz]{Bemerkung}

\begin{document}
	\maketitle
	
	\begin{abstract}
		Wir untersuchen ein numerisches Modell, in dem nach Abzug wachsender glatter Primschalen
		die logarithmischen Beiträge der Restklassen \(1,5,7,11 \pmod{12}\) getrennt betrachtet werden.
		Die bisherigen numerischen Daten legen nahe, dass die normierten mittleren Beiträge der vier
		Klassen nicht zu einem gemeinsamen Grenzwert kollabieren, sondern sich vielmehr eine stabile
		Ordnungsstruktur
		\[
		B>C>A>E
		\]
		ausbildet. Der Artikel formuliert die Definitionen, die zugehörigen Mittelwerte und eine
		präzise numerische Hypothese.
	\end{abstract}
	
	\section{Einleitung}
	
	Für eine natürliche Zahl \(n\) ist es naheliegend, zunächst kleine glatte Faktoren abzutrennen
	und anschließend den verbleibenden Rest nach Restklassen modulo \(12\) zu untersuchen.
	Insbesondere interessiert hier, ob die Klassen
	\[
	1,\ 5,\ 7,\ 11 \pmod{12}
	\]
	asymptotisch gleich behandelt werden oder ob sich nach Schalenreduktion systematische
	Unterschiede zeigen.
	
	Die numerischen Rechnungen deuten darauf hin, dass die Beiträge der vier Klassen im Mittel
	nicht gleich sind. Vielmehr scheint für die normierten logarithmischen Beiträge die Ordnung
	\[
	B>C>A>E
	\]
	stabil zu sein.
	
	\section{Schalen und Restklassenanteile}
	
	\begin{definition}[Primschale]
		Für \(k\in\mathbb N\) sei
		\[
		P_k=\{2,3,5,\dots,p_k\}
		\]
		die Menge der ersten \(k\) Primzahlen.
	\end{definition}
	
	\begin{definition}[Glatter Anteil und Rest]
		Für \(n\in\mathbb N\) definieren wir
		\[
		S_k(n):=\prod_{p\in P_k} p^{v_p(n)}
		\]
		als den \(P_k\)-glatten Anteil und
		\[
		R_k(n):=\frac{n}{S_k(n)}
		\]
		als den schalenreduzierten Rest.
	\end{definition}
	
	\begin{definition}[Restklassenanteile]
		Wir definieren
		\[
		E_k(n):=\prod_{\substack{q^a\parallel R_k(n)\\ q\equiv 1\ (\mathrm{mod}\ 12)}} q^a,
		\]
		\[
		A_k(n):=\prod_{\substack{q^a\parallel R_k(n)\\ q\equiv 5\ (\mathrm{mod}\ 12)}} q^a,
		\]
		\[
		B_k(n):=\prod_{\substack{q^a\parallel R_k(n)\\ q\equiv 7\ (\mathrm{mod}\ 12)}} q^a,
		\]
		\[
		C_k(n):=\prod_{\substack{q^a\parallel R_k(n)\\ q\equiv 11\ (\mathrm{mod}\ 12)}} q^a.
		\]
	\end{definition}
	
	\begin{bemerkung}
		Formal gilt
		\[
		R_k(n)=E_k(n)\,A_k(n)\,B_k(n)\,C_k(n).
		\]
	\end{bemerkung}
	
	\section{Logarithmische Mittelwerte}
	
	\begin{definition}[Grundmenge]
		Für \(N\in\mathbb N\) setzen wir
		\[
		\mathcal M_N:=\{n\le N:\gcd(n,6)=1\}.
		\]
	\end{definition}
	
	\begin{definition}[Logarithmische Mittelwerte]
		Für \(X\in\{E,A,B,C\}\) definieren wir
		\[
		\mu_X^{(k)}(N):=
		\frac{1}{|\mathcal M_N|}
		\sum_{n\in\mathcal M_N}\log X_k(n).
		\]
	\end{definition}
	
	\begin{definition}[Normierte Anteile]
		Wir definieren
		\[
		W_X^{(k)}(N):=
		\frac{\mu_X^{(k)}(N)}
		{\mu_E^{(k)}(N)+\mu_A^{(k)}(N)+\mu_B^{(k)}(N)+\mu_C^{(k)}(N)},
		\qquad X\in\{E,A,B,C\}.
		\]
	\end{definition}
	
	\begin{lemma}
		Für alle \(k,N\) gilt
		\[
		W_E^{(k)}(N)+W_A^{(k)}(N)+W_B^{(k)}(N)+W_C^{(k)}(N)=1.
		\]
	\end{lemma}
	
	\begin{proof}
		Dies folgt unmittelbar aus der Definition.
	\end{proof}
	
	\section{Numerische Hypothese}
	
	\begin{hypothese}[Verschiedene asymptotische Beiträge]
		Für feste oder wachsende Primschalen \(P_k\) existieren asymptotische Grenzwerte
		\[
		\omega_E,\quad \omega_A,\quad \omega_B,\quad \omega_C
		\]
		mit
		\[
		\omega_E+\omega_A+\omega_B+\omega_C=1,
		\]
		so dass für große \(N\) und geeignete Schalen \(k\) gilt
		\[
		W_E^{(k)}(N)\approx \omega_E,\quad
		W_A^{(k)}(N)\approx \omega_A,\quad
		W_B^{(k)}(N)\approx \omega_B,\quad
		W_C^{(k)}(N)\approx \omega_C.
		\]
	\end{hypothese}
	
	\begin{hypothese}[Ordnungsstruktur]
		Die bisherigen numerischen Experimente sprechen für die strikte Ordnung
		\[
		\omega_B>\omega_C>\omega_A>\omega_E.
		\]
	\end{hypothese}
	
	\begin{bemerkung}
		Diese Hypothese ist stärker als die bloße Aussage, dass die Klassen \(A,B,C\) nicht identisch
		sind. Sie postuliert eine stabile Hierarchie der vier Restklassenbeiträge nach Schalenreduktion.
	\end{bemerkung}
	
	\section{Interpretation}
	
	Die Ordnung
	\[
	B>C>A>E
	\]
	bedeutet:
	\begin{itemize}
		\item Die Klasse \(7 \pmod{12}\) trägt im Mittel die größte logarithmische Masse.
		\item Die Klasse \(11 \pmod{12}\) folgt an zweiter Stelle.
		\item Die Klasse \(5 \pmod{12}\) liegt darunter.
		\item Die Klasse \(1 \pmod{12}\), also der energetische Anteil, bleibt am kleinsten.
	\end{itemize}
	
	Dies spricht gegen eine naive Gleichverteilung der vier Klassen nach Schalenreduktion und
	deutet auf eine tieferliegende restklassenspezifische Feinstruktur hin.
	
	\section{Numerisches Testprogramm}
	
	\begin{korollar}
		Zur Prüfung der Hypothesen genügt es, für wachsende \(N\) und wachsende Primschalen \(P_k\)
		die vier Größen
		\[
		\mu_E^{(k)}(N),\quad \mu_A^{(k)}(N),\quad \mu_B^{(k)}(N),\quad \mu_C^{(k)}(N)
		\]
		sowie die normierten Anteile
		\[
		W_E^{(k)}(N),\quad W_A^{(k)}(N),\quad W_B^{(k)}(N),\quad W_C^{(k)}(N)
		\]
		zu berechnen und ihre Ordnung zu verfolgen.
	\end{korollar}
	
	\begin{proof}
		Die Hypothesen betreffen genau diese Größen.
	\end{proof}
	
	\begin{beispiel}
		Falls numerisch für große Schalen wiederholt Werte der Form
		\[
		W_E\approx 0.246,\quad
		W_A\approx 0.248,\quad
		W_B\approx 0.255,\quad
		W_C\approx 0.251
		\]
		auftreten, so spricht dies deutlich für die Hierarchie
		\[
		B>C>A>E.
		\]
	\end{beispiel}
	
	\section{Schluss}
	
	Das hier formulierte Modell trennt eine Zahl zunächst in glatten Anteil und schalenreduzierten
	Rest und misst anschließend die logarithmischen Beiträge der vier Restklassen modulo \(12\).
	Die numerische Evidenz spricht gegen eine volle Symmetrie der Klassen und für eine stabile
	Hierarchie
	\[
	B>C>A>E.
	\]
	
	Ein Beweis dieser Hypothese liegt derzeit nicht vor. Dennoch liefert die Formulierung der
	Mittelwerte und normierten Beiträge eine präzise und überprüfbare Basis für weitere numerische
	und theoretische Untersuchungen.
	\section{Wechsel- und Balanceprimzahlen: strukturelle Vermutungen}
	
	In diesem Abschnitt formulieren wir die aus den numerischen Experimenten gewonnenen
	Arbeitsvermutungen in präziser Form.
	
	\subsection{Grundaufbau}
	
	Sei
	\[
	P_k=\{2,3,5,\dots,p_k\}
	\]
	die \(k\)-te Primschale. Für \(n\in\mathbb N\) definieren wir den Schalenanteil
	\[
	S_k(n):=\prod_{p\in P_k} p^{v_p(n)}
	\]
	und den schalenreduzierten Rest
	\[
	R_k(n):=\frac{n}{S_k(n)}.
	\]
	
	Weiter zerlegen wir den Rest nach den Restklassen modulo \(12\):
	\[
	R_k(n)=E_k(n)\,A_k(n)\,B_k(n)\,C_k(n),
	\]
	wobei
	\[
	E_k(n):=\prod_{\substack{q^a\parallel R_k(n)\\ q\equiv 1\pmod{12}}} q^a,
	\]
	\[
	A_k(n):=\prod_{\substack{q^a\parallel R_k(n)\\ q\equiv 5\pmod{12}}} q^a,
	\quad
	B_k(n):=\prod_{\substack{q^a\parallel R_k(n)\\ q\equiv 7\pmod{12}}} q^a,
	\quad
	C_k(n):=\prod_{\substack{q^a\parallel R_k(n)\\ q\equiv 11\pmod{12}}} q^a.
	\]
	
	Daraus entstehen die drei Vergleichsmodi
	\[
	L_k^{E}(n):=\log E_k(n)-\log S_k(n),
	\]
	\[
	L_k^{ABC}(n):=\log\!\bigl(A_k(n)B_k(n)C_k(n)\bigr)-\log S_k(n),
	\]
	\[
	L_k^{R}(n):=\log R_k(n)-\log S_k(n).
	\]
	
	\subsection{Definitionen}
	
	\begin{definition}[Echte Wechselprimzahl]
		Für einen Modus \(M\in\{E,ABC,R\}\) sei
		\[
		\pi_{\mathrm{wech}}^{M}(n)
		\]
		die erste Schalenprimzahl \(p_k\), bei der das Vorzeichen von \(L_k^{M}(n)\)
		zwischen zwei nichtnulligen Werten wechselt, also
		\[
		\operatorname{sgn}(L_{k-1}^{M}(n))\in\{\pm1\},\qquad
		\operatorname{sgn}(L_k^{M}(n))\in\{\pm1\},
		\]
		und
		\[
		\operatorname{sgn}(L_{k-1}^{M}(n))\neq \operatorname{sgn}(L_k^{M}(n)).
		\]
		Falls kein solcher Wechsel eintritt, setzen wir
		\[
		\pi_{\mathrm{wech}}^{M}(n)=\infty.
		\]
	\end{definition}
	
	\begin{definition}[Nichttriviale Balanceprimzahl]
		Für denselben Modus \(M\) sei
		\[
		\pi_{\mathrm{bal}}^{M}(n)
		\]
		diejenige Schalenprimzahl \(p_k\), für die \(|L_k^{M}(n)|\) minimal ist,
		jedoch nur unter den nichttrivialen Schalenstufen, auf denen beide verglichenen
		Seiten aktiv sind.
		
		Konkret heißt das:
		\begin{itemize}
			\item im \(E\)-Modus:
			\[
			S_k(n)>1,\qquad E_k(n)>1,
			\]
			\item im \(ABC\)-Modus:
			\[
			S_k(n)>1,\qquad A_k(n)B_k(n)C_k(n)>1,
			\]
			\item im \(R\)-Modus:
			\[
			S_k(n)>1,\qquad R_k(n)>1.
			\]
		\end{itemize}
		
		Falls keine solche Stufe existiert, setzen wir
		\[
		\pi_{\mathrm{bal}}^{M}(n)=\infty.
		\]
	\end{definition}
	
	\subsection{Modusreinheit der Restklassen}
	
	\begin{hypothese}[Modusreinheit für Primzahlen]
		Sei \(p>3\) eine Primzahl.
		
		\begin{enumerate}
			\item Falls
			\[
			p\equiv 1\pmod{12},
			\]
			so ist \(p\) ein reiner \(E\)-Träger. In diesem Fall ist numerisch plausibel, dass
			\[
			\pi_{\mathrm{wech}}^{E}(p)=p,
			\qquad
			\pi_{\mathrm{wech}}^{R}(p)=p,
			\qquad
			\pi_{\mathrm{wech}}^{ABC}(p)=\infty,
			\]
			sofern die betrachtete Schalenfolge bis \(p\) reicht.
			
			\item Falls
			\[
			p\equiv 5,7,11\pmod{12},
			\]
			so ist \(p\) ein reiner \(ABC\)-Träger. In diesem Fall ist numerisch plausibel, dass
			\[
			\pi_{\mathrm{wech}}^{ABC}(p)=p,
			\qquad
			\pi_{\mathrm{wech}}^{R}(p)=p,
			\qquad
			\pi_{\mathrm{wech}}^{E}(p)=\infty,
			\]
			sofern die betrachtete Schalenfolge bis \(p\) reicht.
		\end{enumerate}
	\end{hypothese}
	
	\begin{bemerkung}
		Diese Hypothese bringt die numerisch beobachtete Trennung sehr prägnant zum Ausdruck:
		\[
		1\pmod{12}\Rightarrow E\text{-aktiv},
		\qquad
		5,7,11\pmod{12}\Rightarrow ABC\text{-aktiv}.
		\]
		Der \(R\)-Modus vereinigt beide Dynamiken.
	\end{bemerkung}
	
	\subsection{Zweifaktorielle Struktur}
	
	\begin{hypothese}[Zweifaktorzahlen im selben Kanal]
		Sei
		\[
		n=pq
		\]
		mit zwei verschiedenen Primzahlen \(p<q\), und beide Primzahlen seien im selben
		dynamischen Kanal aktiv. Das heißt entweder:
		\begin{itemize}
			\item \(p,q\equiv 1\pmod{12}\) und man betrachtet den \(E\)-Modus,
			\item oder \(p,q\equiv 5,7,11\pmod{12}\) und man betrachtet den \(ABC\)-Modus.
		\end{itemize}
		
		Dann ist numerisch plausibel:
		\[
		\pi_{\mathrm{bal}}^{M}(n)=p,
		\qquad
		\pi_{\mathrm{wech}}^{M}(n)=q.
		\]
	\end{hypothese}
	
	\begin{bemerkung}
		Heuristisch bedeutet dies: Der kleinere Faktor baut die erste nichttriviale
		Schalen-Kern-Konkurrenz auf, während der größere Faktor den eigentlichen Umschlag
		des Dominanzverhältnisses erzwingt.
	\end{bemerkung}
	
	\begin{hypothese}[Gemischte Zweifaktorzahlen]
		Sei
		\[
		n=pq
		\]
		mit
		\[
		p<q,\qquad
		p\equiv 5,7,11\pmod{12},\qquad
		q\equiv 1\pmod{12}.
		\]
		Dann ist im \(R\)-Modus numerisch plausibel:
		\[
		\pi_{\mathrm{bal}}^{R}(n)=p,
		\qquad
		\pi_{\mathrm{wech}}^{R}(n)=q.
		\]
		
		Häufig gilt analog auch im \(E\)-Modus
		\[
		\pi_{\mathrm{bal}}^{E}(n)=p,
		\qquad
		\pi_{\mathrm{wech}}^{E}(n)=q,
		\]
		während im \(ABC\)-Modus der Wechsel bereits bei \(p\) liegt.
	\end{hypothese}
	
	\begin{beispiel}
		Die numerischen Resultate für
		\[
		65=5\cdot 13
		\qquad\text{und}\qquad
		91=7\cdot 13
		\]
		passen sehr gut zu dieser Vermutung:
		\[
		\pi_{\mathrm{bal}}^{R}(65)=5,\qquad \pi_{\mathrm{wech}}^{R}(65)=13,
		\]
		\[
		\pi_{\mathrm{bal}}^{R}(91)=7,\qquad \pi_{\mathrm{wech}}^{R}(91)=13.
		\]
	\end{beispiel}
	
	\subsection{Seltenheit reiner energetischer Balance}
	
	\begin{hypothese}[Seltenheit nichttrivialer Balance im \(E\)-Modus]
		Für viele Zahlen \(n\) existiert keine nichttriviale Balanceprimzahl im reinen
		\(E\)-Modus, also
		\[
		\pi_{\mathrm{bal}}^{E}(n)=\infty.
		\]
		Dagegen existiert im \(R\)-Modus deutlich häufiger eine nichttriviale Balance.
	\end{hypothese}
	
	\begin{bemerkung}
		Dies legt nahe, dass der reine energetische Anteil oft zu schmal ist, um gegen die
		Schale eine robuste Balance auszubilden. Der volle Rest \(R_k(n)\) erfasst die
		Kern-Schalen-Dynamik wesentlich zuverlässiger.
	\end{bemerkung}
	
	\subsection{Hauptvermutung}
	
	\begin{hypothese}[Balance unten, Wechsel oben]
		Die bisherigen numerischen Daten deuten auf folgende allgemeine Struktur hin:
		
		Für eine Zahl \(n\) mit aktiven Primfaktoren in einem gegebenen Modus \(M\) gilt
		typischerweise:
		\[
		\pi_{\mathrm{bal}}^{M}(n)\approx \text{kleinster aktiver Primfaktor},
		\qquad
		\pi_{\mathrm{wech}}^{M}(n)\approx \text{größter aktiver Primfaktor}.
		\]
		
		Insbesondere scheint bei zweifaktoriellen oder schwach zusammengesetzten Zahlen
		eine systematische Trennung zwischen Balance- und Wechselprimzahl aufzutreten:
		die Balance liegt unten, der eigentliche Wechsel oben.
	\end{hypothese}
	
	\begin{bemerkung}
		Diese Vermutung bündelt die numerisch sichtbare geometrische Kernaussage des Modells:
		Die Schale wird zuerst durch kleine Faktoren aufgebaut, während größere Faktoren
		den Kern länger stabil halten und deshalb den eigentlichen Umschlag bestimmen.
	\end{bemerkung}
	
	\section{Der semiprime Fall: Balance unten, Wechsel oben}
	
	In diesem Abschnitt formulieren wir die numerisch vollständig bestätigte Struktur
	für semiprimale Zahlen als präzise Vermutung mit Beweisskizze.
	
	\begin{definition}[Aktiver Modus für semiprimale Zahlen]
		Sei
		\[
		n=pq,\qquad p<q,
		\]
		das Produkt zweier verschiedener Primzahlen \(>3\).
		
		Wir nennen einen Modus \(M\in\{E,ABC,R\}\) \emph{aktiv}, wenn beide Faktoren von \(n\)
		im betreffenden Vergleich tatsächlich Masse tragen. Konkret:
		
		\begin{itemize}
			\item Für \(p,q\equiv 1\pmod{12}\) ist der \(E\)-Modus aktiv.
			\item Für \(p,q\equiv 5,7,11\pmod{12}\) ist der \(ABC\)-Modus aktiv.
			\item Für gemischte Produkte mit einem Faktor in \(ABC\) und einem Faktor in \(E\)
			ist der \(R\)-Modus stets aktiv.
		\end{itemize}
	\end{definition}
	
	\begin{satz}[Semiprime Klein-Groß-Regel]
		\label{satz:semiprime_klein_gross}
		Sei
		\[
		n=pq,\qquad p<q,
		\]
		das Produkt zweier verschiedener Primzahlen \(>3\), und sei \(M\in\{E,ABC,R\}\)
		der zu \(n\) gehörige aktive Modus. Dann gilt numerisch konsistent:
		
		\[
		\pi_{\mathrm{bal}}^{M}(n)=p,
		\qquad
		\pi_{\mathrm{wech}}^{M}(n)=q.
		\]
		
		Das heißt:
		\begin{enumerate}
			\item Die nichttriviale Balanceprimzahl liegt beim kleineren aktiven Primfaktor.
			\item Die echte Wechselprimzahl liegt beim größeren aktiven Primfaktor.
		\end{enumerate}
	\end{satz}
	
	\begin{bemerkung}
		Die bisherigen semiprimen Tests zeigen diese Regel in allen getesteten Fällen:
		\begin{itemize}
			\item \(E\cdot E\) im \(E\)-Modus: \(100\%\),
			\item \(ABC\cdot ABC\) im \(ABC\)-Modus: \(100\%\),
			\item \(ABC\cdot E\) im \(R\)-Modus: \(100\%\),
			\item \(ABC\cdot E\) im \(E\)-Modus ebenfalls \(100\%\) für Balance unten und Wechsel oben.
		\end{itemize}
	\end{bemerkung}
	
	\begin{proof}[Beweisskizze]
		Wir geben eine heuristisch-elementare Begründung.
		
		\medskip
		\noindent
		\textbf{1. Aufbau der Schale.}
		Für
		\[
		n=pq,\qquad p<q,
		\]
		wird bei wachsender Primschale
		\[
		P_k=\{2,3,5,\dots,p_k\}
		\]
		zunächst kein aktiver Faktor erfasst, solange \(p_k<p\). In diesem Bereich bleibt der
		aktive Rest vollständig erhalten.
		
		Sobald jedoch die Schale den kleineren Faktor \(p\) erreicht, gilt:
		\[
		S_k(n)=p,\qquad R_k(n)=q.
		\]
		Damit tritt erstmals eine echte Konkurrenz zwischen Schale und aktivem Rest auf.
		
		\medskip
		\noindent
		\textbf{2. Warum die Balance unten liegt.}
		Vor dem Eintritt von \(p\) in die Schale ist die Schale zu klein, um mit dem aktiven Rest
		zu konkurrieren; nach Eintritt von \(p\) ist die Schale erstmals nichttrivial, aber der
		größere Faktor \(q\) verbleibt noch vollständig im Rest. Daher ist genau in diesem Bereich
		die beste Balance zu erwarten.
		
		Im semiprimen Fall gibt es nach Erreichen von \(p\) nur noch einen aktiven Restfaktor,
		nämlich \(q\). Deshalb ist die erste nichttriviale Stufe zugleich die natürlichste Stelle,
		an der
		\[
		|L_k^M(n)|
		\]
		minimal wird. Daraus folgt heuristisch
		\[
		\pi_{\mathrm{bal}}^{M}(n)=p.
		\]
		
		\medskip
		\noindent
		\textbf{3. Warum der Wechsel oben liegt.}
		Solange \(q\) noch nicht in die Schale eingetreten ist, bleibt im aktiven Modus eine
		nichttriviale Restmasse erhalten. Erst beim Eintritt von \(q\) wird auch dieser letzte
		aktive Faktor aus dem Rest entfernt bzw. durch die Schale absorbiert. Genau dort kippt das
		Dominanzverhältnis endgültig.
		
		Daher ist zu erwarten, dass der erste echte Vorzeichenwechsel zwischen zwei nichtnulligen
		Werten an der Schalenprimzahl \(q\) eintritt:
		\[
		\pi_{\mathrm{wech}}^{M}(n)=q.
		\]
		
		\medskip
		\noindent
		\textbf{4. Modusabhängigkeit.}
		Welche der Größen
		\[
		L_k^E(n),\qquad L_k^{ABC}(n),\qquad L_k^R(n)
		\]
		tatsächlich diese Dynamik tragen, hängt allein davon ab, in welcher modulo-\(12\)-Klasse
		die Faktoren \(p,q\) liegen:
		
		\begin{itemize}
			\item Bei \(E\cdot E\) liegt die gesamte aktive Masse im \(E\)-Kanal.
			\item Bei \(ABC\cdot ABC\) liegt sie im \(ABC\)-Kanal.
			\item Bei gemischten Zahlen \(ABC\cdot E\) ist der \(R\)-Modus der natürliche Gesamtmodus,
			weil er beide Beiträge vereinigt.
		\end{itemize}
		
		Damit ergibt sich die Klein-Groß-Regel im jeweils aktiven Modus.
	\end{proof}
	
	\begin{korollar}[Modusreine semiprimale Fälle]
		\label{kor:modusrein_semiprime}
		Für semiprimale Zahlen \(n=pq\) mit \(p<q\) gilt:
		
		\begin{enumerate}
			\item Falls
			\[
			p,q\equiv 1\pmod{12},
			\]
			dann ist der \(E\)-Modus aktiv und
			\[
			\pi_{\mathrm{bal}}^{E}(n)=p,\qquad
			\pi_{\mathrm{wech}}^{E}(n)=q.
			\]
			
			\item Falls
			\[
			p,q\equiv 5,7,11\pmod{12},
			\]
			dann ist der \(ABC\)-Modus aktiv und
			\[
			\pi_{\mathrm{bal}}^{ABC}(n)=p,\qquad
			\pi_{\mathrm{wech}}^{ABC}(n)=q.
			\]
			
			\item Falls
			\[
			p\equiv 5,7,11\pmod{12},\qquad q\equiv 1\pmod{12},
			\]
			dann ist der \(R\)-Modus aktiv und
			\[
			\pi_{\mathrm{bal}}^{R}(n)=p,\qquad
			\pi_{\mathrm{wech}}^{R}(n)=q.
			\]
		\end{enumerate}
	\end{korollar}
	
	\begin{beispiel}
		Die numerischen Beispiele
		\[
		35=5\cdot 7,\qquad
		55=5\cdot 11,\qquad
		77=7\cdot 11
		\]
		bestätigen im \(ABC\)-Modus:
		\[
		\pi_{\mathrm{bal}}^{ABC}(35)=5,\quad \pi_{\mathrm{wech}}^{ABC}(35)=7,
		\]
		\[
		\pi_{\mathrm{bal}}^{ABC}(55)=5,\quad \pi_{\mathrm{wech}}^{ABC}(55)=11,
		\]
		\[
		\pi_{\mathrm{bal}}^{ABC}(77)=7,\quad \pi_{\mathrm{wech}}^{ABC}(77)=11.
		\]
		
		Ebenso bestätigen die gemischten Beispiele
		\[
		65=5\cdot 13,\qquad 91=7\cdot 13
		\]
		im \(R\)-Modus:
		\[
		\pi_{\mathrm{bal}}^{R}(65)=5,\quad \pi_{\mathrm{wech}}^{R}(65)=13,
		\]
		\[
		\pi_{\mathrm{bal}}^{R}(91)=7,\quad \pi_{\mathrm{wech}}^{R}(91)=13.
		\]
	\end{beispiel}
	
	\begin{bemerkung}
		Ein vollständiger strenger Beweis müsste die genaue Monotonie oder wenigstens die
		Eindeutigkeit des Minimums von
		\[
		|L_k^M(n)|
		\]
		zwischen den Schalenstufen \(p\) und \(q\) kontrollieren. Die bisherige numerische Evidenz
		legt jedoch nahe, dass der semiprime Fall wesentlich rigider ist als der allgemeine
		mehrfaktorielle Fall.
	\end{bemerkung}
	
	\subsection{Der Primzahlfall}
	
	\begin{satz}[Primzahlen als modusreine Elementarfälle]
		\label{prop:primzahlfall}
		Sei \(p>3\) eine Primzahl.
		
		\begin{enumerate}
			\item Falls
			\[
			p\equiv 1\pmod{12},
			\]
			so ist \(p\) im dynamischen Sinn ein reiner \(E\)-Träger. In diesem Fall gilt
			im betrachteten Schalenmodell numerisch:
			\[
			\pi_{\mathrm{wech}}^{E}(p)=p,\qquad
			\pi_{\mathrm{wech}}^{R}(p)=p,\qquad
			\pi_{\mathrm{wech}}^{ABC}(p)=\infty.
			\]
			
			\item Falls
			\[
			p\equiv 5,7,11\pmod{12},
			\]
			so ist \(p\) ein reiner \(ABC\)-Träger. In diesem Fall gilt numerisch:
			\[
			\pi_{\mathrm{wech}}^{ABC}(p)=p,\qquad
			\pi_{\mathrm{wech}}^{R}(p)=p,\qquad
			\pi_{\mathrm{wech}}^{E}(p)=\infty.
			\]
		\end{enumerate}
	\end{satz}
	
	\begin{proof}[Beweisskizze]
		Wir betrachten eine Primzahl \(p>3\) entlang der Schalenfolge
		\[
		P_k=\{2,3,5,\dots,p_k\}.
		\]
		
		Solange \(p_k<p\), gilt stets
		\[
		S_k(p)=1,\qquad R_k(p)=p.
		\]
		Die gesamte aktive Masse liegt also vollständig im Rest.
		
		\medskip
		\noindent
		\textbf{Fall 1: \(p\equiv 1\pmod{12}\).}
		Dann ist
		\[
		E_k(p)=p,\qquad A_k(p)=B_k(p)=C_k(p)=1
		\]
		für alle Schalenstufen mit \(p_k<p\). Daher ist der \(E\)-Modus aktiv, während
		der \(ABC\)-Modus trivial bleibt. Sobald die Schale die Primzahl \(p\) selbst erreicht,
		wird \(p\) aus dem Rest entfernt, also
		\[
		R_k(p)=1.
		\]
		Damit tritt der eigentliche Umschlag im \(E\)-Modus und ebenso im \(R\)-Modus
		genau bei der Schalenprimzahl \(p\) auf:
		\[
		\pi_{\mathrm{wech}}^{E}(p)=p,\qquad \pi_{\mathrm{wech}}^{R}(p)=p.
		\]
		Da im \(ABC\)-Modus keine aktive Masse vorliegt, bleibt dort
		\[
		\pi_{\mathrm{wech}}^{ABC}(p)=\infty.
		\]
		
		\medskip
		\noindent
		\textbf{Fall 2: \(p\equiv 5,7,11\pmod{12}\).}
		Dann liegt die gesamte Restmasse im \(ABC\)-Kanal, also
		\[
		A_k(p)B_k(p)C_k(p)=p,\qquad E_k(p)=1
		\]
		für alle Schalenstufen mit \(p_k<p\). Daher ist jetzt der \(ABC\)-Modus aktiv,
		während der \(E\)-Modus trivial bleibt. Sobald die Schale \(p\) selbst erreicht,
		verschwindet auch hier der Rest, und der Wechsel tritt genau bei \(p\) auf:
		\[
		\pi_{\mathrm{wech}}^{ABC}(p)=p,\qquad \pi_{\mathrm{wech}}^{R}(p)=p.
		\]
		Dagegen bleibt
		\[
		\pi_{\mathrm{wech}}^{E}(p)=\infty.
		\]
		
		In beiden Fällen ist also die Primzahl selbst genau die erste Schalenprimzahl,
		an der der jeweils aktive Modus seinen Wechsel zeigt.
	\end{proof}
	
	\begin{korollar}[Modusklassifikation der Primzahlen]
		\label{kor:modusklassifikation_primzahlen}
		Für Primzahlen \(p>3\) gilt:
		\[
		p\equiv 1\pmod{12}
		\quad\Longrightarrow\quad
		\text{\(p\) ist \(E\)-aktiv und \(ABC\)-inaktiv,}
		\]
		\[
		p\equiv 5,7,11\pmod{12}
		\quad\Longrightarrow\quad
		\text{\(p\) ist \(ABC\)-aktiv und \(E\)-inaktiv.}
		\]
		Der \(R\)-Modus ist in beiden Fällen aktiv und vereinigt die jeweilige Dynamik.
	\end{korollar}
	
	\begin{bemerkung}
		Der Primzahlfall liefert die elementaren Bausteine des gesamten Modells. Die
		semiprime Klein-Groß-Regel aus Satz~\ref{satz:semiprime_klein_gross} kann als
		erste nichttriviale Erweiterung dieses modusreinen Primzahlverhaltens verstanden
		werden.
	\end{bemerkung}
	
	\subsection{Der Fall \(p^a q^b\): Balance unten, Wechsel nach Log-Masse}
	
	Die semiprimale Klein-Groß-Regel
	\[
	\pi_{\mathrm{bal}}=p,\qquad \pi_{\mathrm{wech}}=q
	\]
	für
	\[
	n=pq,\qquad p<q,
	\]
	erweist sich bei höheren Exponenten als nur teilweise stabil. Die numerischen Tests für
	\[
	n=p^a q^b
	\]
	zeigen jedoch eine deutlich verfeinerte Struktur.
	
	\begin{hypothese}[Robuste Balanceprimzahl im Zweifaktorfall]
		\label{hyp:paqb_balance}
		Sei
		\[
		n=p^a q^b,\qquad p<q,
		\]
		mit Primzahlen \(p,q>3\) und Exponenten \(a,b\ge 1\). Sei \(M\in\{E,ABC,R\}\) der
		zu \(n\) gehörige aktive Modus.
		
		Dann ist numerisch stark plausibel, dass die nichttriviale Balanceprimzahl weiterhin durch
		den kleineren Primfaktor gegeben ist:
		\[
		\pi_{\mathrm{bal}}^{M}(n)=p.
		\]
	\end{hypothese}
	
	\begin{bemerkung}
		Die numerischen Tests für Zahlen der Form \(p^a q^b\) zeigen in den aktiven Modi für die
		Familien \(E\cdot E\), \(ABC\cdot ABC\) und \(ABC\cdot E\) eine volle Stabilität der
		Balance-Regel:
		\[
		\pi_{\mathrm{bal}}^{M}(n)=p.
		\]
		Insbesondere ist die Balanceprimzahl wesentlich robuster als die Wechselprimzahl.
	\end{bemerkung}
	
	\begin{hypothese}[Log-Massen-Kriterium für die Wechselprimzahl]
		\label{hyp:paqb_wechsel}
		Sei
		\[
		n=p^a q^b,\qquad p<q,
		\]
		mit Primzahlen \(p,q>3\), Exponenten \(a,b\ge 1\), und sei \(M\) der aktive Modus.
		
		Dann wird die Wechselprimzahl nicht mehr rein durch die Ordnung \(p<q\), sondern durch die
		gewichteten logarithmischen Beiträge
		\[
		a\log p
		\qquad\text{und}\qquad
		b\log q
		\]
		bestimmt. Numerisch legt dies folgende Regel nahe:
		\[
		\pi_{\mathrm{wech}}^{M}(n)=
		\begin{cases}
			p, & \text{falls } a\log p \ge b\log q,\\[1mm]
			q, & \text{falls } a\log p < b\log q.
		\end{cases}
		\]
	\end{hypothese}
	
	\begin{bemerkung}
		Diese Regel verallgemeinert die semiprimale Situation \(a=b=1\). In diesem Spezialfall gilt
		\[
		\log p<\log q,
		\]
		also automatisch
		\[
		\pi_{\mathrm{wech}}^{M}(pq)=q.
		\]
		Damit wird die semiprimale Klein-Groß-Regel als Spezialfall des Log-Massen-Kriteriums
		sichtbar.
	\end{bemerkung}
	
	\begin{proof}[Heuristische Begründung]
		Wir skizzieren die zugrunde liegende Dynamik.
		
		\medskip
		\noindent
		\textbf{1. Erste Konkurrenzlage.}
		Solange die Schale den kleineren Primfaktor \(p\) noch nicht enthält, bleibt der aktive
		Rest vollständig erhalten. Beim Eintritt von \(p\) in die Schale wird erstmals ein
		nichttrivialer Schalenanteil erzeugt, während der Rest noch mindestens den Beitrag von
		\(q^b\) enthält. Daher ist der Eintritt von \(p\) die natürlichste Stelle für die erste
		echte Balance zwischen Schale und Kern.
		
		Dies erklärt heuristisch die Stabilität von
		\[
		\pi_{\mathrm{bal}}^{M}(n)=p.
		\]
		
		\medskip
		\noindent
		\textbf{2. Dominanz nach Log-Masse.}
		Nach Eintritt von \(p\) ist die Schalenmasse bereits von der Größe
		\[
		a\log p,
		\]
		während im Rest mindestens noch die Masse
		\[
		b\log q
		\]
		verbleibt. Ist
		\[
		a\log p \ge b\log q,
		\]
		so ist bereits beim Eintritt von \(p\) die Schale mindestens so stark wie der verbleibende
		aktive Rest, und der Wechsel tritt sofort bei \(p\) ein.
		
		Ist dagegen
		\[
		a\log p < b\log q,
		\]
		so hält der Restfaktor \(q^b\) die Dominanz noch aufrecht, und der eigentliche Wechsel wird
		erst beim Eintritt von \(q\) in die Schale erzwungen.
		
		\medskip
		\noindent
		\textbf{3. Interpretation.}
		Die Balance wird also geometrisch durch den ersten relevanten Eintritt in die Schale
		bestimmt, der Wechsel dagegen durch die relative logarithmische Masse der aktiven Faktoren.
	\end{proof}
	
	\begin{korollar}[Semiprime als Spezialfall]
		\label{kor:semiprime_spezialfall_logmasse}
		Für
		\[
		n=pq,\qquad p<q,
		\]
		gilt \(a=b=1\). Damit folgt aus Hypothese~\ref{hyp:paqb_wechsel}
		\[
		\log p<\log q,
		\]
		und daher
		\[
		\pi_{\mathrm{wech}}^{M}(pq)=q.
		\]
		Zusammen mit Hypothese~\ref{hyp:paqb_balance} ergibt sich wieder
		\[
		\pi_{\mathrm{bal}}^{M}(pq)=p,\qquad
		\pi_{\mathrm{wech}}^{M}(pq)=q.
		\]
	\end{korollar}
	
	\begin{beispiel}
		Die numerischen Gegenbeispiele zur naiven Regel
		\[
		\pi_{\mathrm{wech}}^{M}(n)=q
		\]
		werden gerade durch hohe Exponenten des kleineren Faktors erklärt.
		
		Zum Beispiel:
		\[
		n=5^2\cdot 7.
		\]
		Hier gilt
		\[
		2\log 5 > \log 7,
		\]
		und numerisch tritt der Wechsel bereits bei \(5\) auf:
		\[
		\pi_{\mathrm{wech}}^{M}(n)=5.
		\]
		
		Ebenso für
		\[
		n=5^3\cdot 11^2,
		\]
		denn
		\[
		3\log 5 > 2\log 11.
		\]
		Auch hier wird numerisch
		\[
		\pi_{\mathrm{wech}}^{M}(n)=5
		\]
		beobachtet.
	\end{beispiel}
	
	\begin{bemerkung}
		Die bisherige Evidenz legt nahe, dass die Balanceprimzahl der eigentliche robuste Invariant
		des Zweifaktormodells ist, während die Wechselprimzahl das feinere, exponentenempfindliche
		Objekt darstellt.
	\end{bemerkung}
	
	Die numerischen Experimente zeigen, dass sich das Schalenwechsel-Modell bereits im
	Zweifaktorfall überraschend klar strukturiert. Für semiprimale Zahlen
	\[
	n=pq,\qquad p<q,
	\]
	tritt im jeweils aktiven Modus eine saubere Trennung zwischen Balance- und
	Wechselprimzahl auf: Die nichttriviale Balance liegt beim kleineren Primfaktor,
	der eigentliche Wechsel beim größeren. Für Zahlen der Form
	\[
	n=p^a q^b
	\]
	bleibt diese Regel nur teilweise erhalten. Während die Balanceprimzahl numerisch
	sehr stabil beim kleineren Primfaktor verbleibt, reagiert die Wechselprimzahl
	empfindlich auf die Exponenten. Die Daten legen nahe, dass der Wechsel nicht
	mehr allein durch die Ordnung \(p<q\), sondern durch die gewichteten logarithmischen
	Beiträge
	\[
	a\log p
	\qquad\text{und}\qquad
	b\log q
	\]
	gesteuert wird. Damit erscheint die Balanceprimzahl als robuster geometrischer
	Invariant der ersten echten Schalen-Kern-Konkurrenz, während die Wechselprimzahl
	das feinere dynamische Objekt darstellt, das auf die Massenverteilung der
	aktiven Faktoren reagiert.
	
	Die bisherigen numerischen Resultate sprechen dafür, dass das Schalenwechsel-Modell
	zwei verschiedene Ebenen arithmetischer Struktur sichtbar macht. Erstens gibt es
	mit der Balanceprimzahl einen bemerkenswert stabilen Invariant: Im Zweifaktorfall
	liegt die erste nichttriviale Balance fast durchgängig beim kleineren aktiven
	Primfaktor. Zweitens zeigt die Wechselprimzahl eine deutlich feinere Dynamik.
	Im semiprimen Fall fällt sie mit dem größeren Primfaktor zusammen; bei höheren
	Exponenten wird sie jedoch durch die logarithmische Massenverteilung der Faktoren
	beeinflusst. Dadurch entsteht eine natürliche Trennung zwischen geometrischer
	Eintrittsstruktur und gewichteter Dominanzstruktur. Gerade diese Zweiteilung
	dürfte für die weitere Entwicklung des Modells zentral sein.
	
	\subsection{Der Fall \(p^a q^b\): robuste Balance und log-massen-gesteuerter Wechsel}
	
	Die numerischen Tests für Zahlen der Form
	\[
	n=p^a q^b,\qquad p<q,
	\]
	zeigen im untersuchten Parameterbereich eine klare Zweiteilung zwischen Balance- und
	Wechselprimzahl. Während die Balanceprimzahl bemerkenswert stabil beim kleineren
	Primfaktor verbleibt, wird die Wechselprimzahl durch die gewichteten logarithmischen
	Beiträge der beiden Faktoren bestimmt.
	
	\begin{satz}[Numerisch bestätigte Balance-Regel für \(p^a q^b\)]
		\label{satz:paqb_balance}
		Sei
		\[
		n=p^a q^b,\qquad p<q,
		\]
		mit Primzahlen \(p,q>3\) und Exponenten \(a,b\ge 1\). Sei \(M\in\{E,ABC,R\}\) der
		zu \(n\) gehörige aktive Modus.
		
		Dann gilt in allen bisher getesteten Fällen:
		\[
		\pi_{\mathrm{bal}}^{M}(n)=p.
		\]
		Mit anderen Worten: Die nichttriviale Balanceprimzahl wird durchgehend vom kleineren
		Primfaktor bestimmt.
	\end{satz}
	
	\begin{proof}[Numerische Begründung]
		Die systematischen Tests über die Familien
		\[
		E\cdot E,\qquad ABC\cdot ABC,\qquad ABC\cdot E
		\]
		zeigen in den jeweiligen aktiven Modi für alle untersuchten Zahlen der Form
		\[
		p^a q^b
		\]
		eine volle Trefferquote der Regel
		\[
		\pi_{\mathrm{bal}}^{M}(n)=p.
		\]
		Dies gilt sowohl in den reinen Kanälen \(E\) bzw. \(ABC\) als auch im Gesamtmodus \(R\).
	\end{proof}
	
	\begin{satz}[Numerisch bestätigte Log-Massen-Regel für die Wechselprimzahl]
		\label{satz:paqb_logmasse}
		Unter denselben Voraussetzungen wie in Satz~\ref{satz:paqb_balance} gilt im aktiven Modus
		\(M\) numerisch:
		
		\[
		\pi_{\mathrm{wech}}^{M}(n)=
		\begin{cases}
			p, & \text{falls } a\log p \ge b\log q,\\[1mm]
			q, & \text{falls } a\log p < b\log q.
		\end{cases}
		\]
		
		Die Wechselprimzahl wird also nicht mehr nur durch die Ordnung \(p<q\), sondern durch die
		gewichteten logarithmischen Massen
		\[
		a\log p
		\qquad\text{und}\qquad
		b\log q
		\]
		bestimmt.
	\end{satz}
	
	\begin{proof}[Numerische Begründung]
		Die naive semiprimale Regel
		\[
		\pi_{\mathrm{wech}}^{M}(n)=q
		\]
		scheitert bei höheren Exponenten systematisch. Die isolierten Gegenbeispiele treten genau
		in den Fällen auf, in denen der kleinere Primfaktor durch seinen Exponenten genügend
		logarithmische Masse trägt, um bereits beim Eintritt in die Schale dominant zu werden.
		
		Für die verfeinerte Vorhersage
		\[
		\pi_{\mathrm{wech}}^{M}(n)=
		\begin{cases}
			p,& a\log p \ge b\log q,\\
			q,& a\log p < b\log q
		\end{cases}
		\]
		wurden im getesteten Bereich keine Gegenbeispiele gefunden. Die Trefferquote der
		Log-Massen-Regel beträgt in allen ausgewerteten Familien und in beiden Bereichen
		\[
		\frac{a\log p}{b\log q}\ge 1
		\qquad\text{bzw.}\qquad
		\frac{a\log p}{b\log q}<1
		\]
		jeweils \(100\%\).
	\end{proof}
	
	\begin{korollar}[Der semiprimale Fall als Spezialfall]
		\label{kor:paqb_semiprime}
		Für
		\[
		n=pq,\qquad p<q,
		\]
		gilt \(a=b=1\). Daher ist
		\[
		\log p<\log q,
		\]
		also folgt aus Satz~\ref{satz:paqb_logmasse}
		\[
		\pi_{\mathrm{wech}}^{M}(pq)=q.
		\]
		Zusammen mit Satz~\ref{satz:paqb_balance} erhält man erneut
		\[
		\pi_{\mathrm{bal}}^{M}(pq)=p,\qquad
		\pi_{\mathrm{wech}}^{M}(pq)=q.
		\]
		Der semiprimale Fall erscheint somit als Spezialfall der allgemeinen \(p^a q^b\)-Struktur.
	\end{korollar}
	
	\begin{beispiel}
		Die klassischen Gegenbeispiele zur naiven Regel \(\pi_{\mathrm{wech}}=q\) werden durch das
		Log-Massen-Kriterium exakt erklärt.
		
		\begin{enumerate}
			\item Für
			\[
			n=5^2\cdot 7
			\]
			gilt
			\[
			2\log 5 > \log 7.
			\]
			Daher sagt Satz~\ref{satz:paqb_logmasse} voraus:
			\[
			\pi_{\mathrm{wech}}(n)=5,
			\]
			was numerisch bestätigt wird.
			
			\item Für
			\[
			n=5^3\cdot 11^2
			\]
			gilt
			\[
			3\log 5 > 2\log 11.
			\]
			Also folgt wiederum
			\[
			\pi_{\mathrm{wech}}(n)=5,
			\]
			ebenfalls in Übereinstimmung mit den numerischen Daten.
			
			\item Für
			\[
			n=7^2\cdot 11
			\]
			gilt
			\[
			2\log 7 > \log 11,
			\]
			also tritt der Wechsel bereits bei \(7\) auf:
			\[
			\pi_{\mathrm{wech}}(n)=7.
			\]
		\end{enumerate}
	\end{beispiel}
	
	\begin{bemerkung}
		Die numerischen Resultate legen nahe, dass die Balanceprimzahl der robuste geometrische
		Invariant des Zweifaktormodells ist, während die Wechselprimzahl das feinere dynamische
		Objekt darstellt, das auf die logarithmische Massenverteilung der aktiven Faktoren
		reagiert.
	\end{bemerkung}
	
	\begin{bemerkung}
		Ein strenger Beweis von Satz~\ref{satz:paqb_logmasse} müsste zeigen, dass der erste echte
		Vorzeichenwechsel des jeweiligen Modus tatsächlich genau an der Primzahl eintritt, deren
		gewichteter logarithmischer Beitrag den dominanten Restblock kontrolliert. Die bisherige
		Datenlage spricht jedoch sehr stark dafür, dass dies im Zweifaktorfall die richtige
		Struktur ist.
	\end{bemerkung}
	
	\subsection{Übergang zum Dreifaktorfall}
	
	Nach dem Primzahlfall und dem vollständig durchgerechneten Zweifaktorfall
	\[
	n=p^a q^b
	\]
	stellt der Dreifaktorfall die nächste natürliche Stufe der Analyse dar. Die bisherigen
	numerischen Resultate legen nahe, dass die Balanceprimzahl auch dort weiterhin vom
	kleinsten aktiven Primfaktor kontrolliert werden dürfte. Offen ist jedoch, wie sich die
	Wechselprimzahl in Anwesenheit eines dritten Faktors verhält. Im Zweifaktorfall wird sie
	durch die gewichteten logarithmischen Beiträge
	\[
	a\log p
	\qquad\text{und}\qquad
	b\log q
	\]
	bestimmt; im Dreifaktorfall ist daher zu erwarten, dass statt eines bloßen
	Zweiervergleichs eine Konkurrenz mehrerer aktiver Restblöcke auftritt. Insbesondere stellt
	sich die Frage, ob der eigentliche Wechsel weiterhin durch den größten Primfaktor
	hervorgerufen wird, oder ob vielmehr der größte verbleibende logarithmische Massenblock
	entscheidend ist. Der Übergang von \(p^a q^b\) zu
	\[
	p^a q^b r^c
	\]
	markiert damit den Punkt, an dem aus einer zweistufigen Schalen-Kern-Dynamik eine echte
	mehrstufige Dominanzstruktur wird.
	
	\subsection{Übergang zum Dreifaktorfall}
	
	Der Zweifaktorfall
	\[
	n=p^a q^b
	\]
	zeigt bereits eine bemerkenswert klare Trennung: Die Balanceprimzahl ist geometrisch stabil
	und liegt beim kleineren Primfaktor, während die Wechselprimzahl von der logarithmischen
	Massenverteilung der aktiven Faktoren abhängt. Damit ist der nächste Schritt nahezu
	zwingend: der Übergang zu Zahlen der Form
	\[
	n=p^a q^b r^c.
	\]
	Hier ist nicht länger nur zu entscheiden, welcher von zwei Faktoren den Rest dominiert.
	Vielmehr entsteht eine dreifache Konkurrenzlage, in der Balance, Vorzeichenwechsel und
	Restdominanz möglicherweise auf verschiedene Primfaktoren verteilt werden. Gerade dieser
	Dreifaktorfall dürfte zeigen, ob die bisher beobachtete Zweiteilung zwischen geometrischem
	Eintrittspunkt und gewichteter Dominanz ein Spezialphänomen des Zweifaktormodells ist oder
	den Anfang einer allgemeineren Schalenstruktur bildet.
	
	\subsection{Erste Vermutungen zum Dreifaktorfall}
	
	Die numerischen Resultate für Zahlen der Form
	\[
	n=p^a q^b r^c,\qquad p<q<r,
	\]
	zeigen, dass die Zweifaktorregeln nicht direkt auf den Dreifaktorfall übergehen. Weder
	liegt die Balanceprimzahl allgemein beim kleinsten Primfaktor \(p\), noch wird die
	Wechselprimzahl zuverlässig durch den größten Primfaktor \(r\) oder allein durch die
	größte logarithmische Einzelmasse bestimmt. Stattdessen deutet sich eine mehrstufige
	Dynamik an, in der verschiedene Primfaktoren verschiedene Rollen übernehmen können.
	
	\begin{hypothese}[Mittlere Balance im Dreifaktorfall]
		\label{hyp:dreifaktor_balance_mitte}
		Sei
		\[
		n=p^a q^b r^c,\qquad p<q<r,
		\]
		mit Primzahlen \(p,q,r>3\), und sei \(M\in\{E,ABC,R\}\) der aktive Modus.
		
		Dann ist numerisch plausibel, dass die nichttriviale Balanceprimzahl im Dreifaktorfall
		häufig nicht mehr durch den kleinsten Faktor \(p\), sondern durch einen \emph{mittleren}
		aktiven Faktor bestimmt wird. Im einfachsten Fall
		\[
		n=pqr
		\]
		spricht die bisherige Evidenz insbesondere dafür, dass oft
		\[
		\pi_{\mathrm{bal}}^{M}(n)=q
		\]
		gilt.
	\end{hypothese}
	
	\begin{bemerkung}
		Die bisherige Datenlage legt nahe, dass der Dreifaktorfall nicht länger nur eine
		Konkurrenz zwischen früher Schalenbildung und spätem Restumschlag darstellt. Vielmehr
		scheint sich eine Zwischenlage auszubilden, in der der mittlere Faktor die beste Balance
		zwischen bereits aufgebauter Schale und verbleibender Restmasse realisiert.
	\end{bemerkung}
	
	\begin{hypothese}[Getrennte Rollen von Eintritt, Balance und Wechsel]
		\label{hyp:dreifaktor_rollentrennung}
		Für Zahlen der Form
		\[
		n=p^a q^b r^c,\qquad p<q<r,
		\]
		treten im aktiven Modus \(M\) typischerweise drei verschiedene dynamische Rollen auf:
		
		\begin{enumerate}
			\item ein \emph{Eintrittsfaktor}, bei dessen Erreichen die erste echte Konkurrenz zwischen
			Schale und Rest beginnt,
			\item ein \emph{Balancefaktor}, bei dem \(|L_k^M(n)|\) minimal wird,
			\item ein \emph{Wechselfaktor}, bei dem der erste echte Vorzeichenwechsel eintritt.
		\end{enumerate}
		
		Im Dreifaktorfall müssen diese drei Rollen im Allgemeinen nicht zusammenfallen.
		Insbesondere ist numerisch plausibel, dass der Balancefaktor häufig der mittlere Faktor
		\(q\) ist, während der Wechselfaktor je nach Massenverteilung zwischen \(q\) und \(r\)
		variieren kann.
	\end{hypothese}
	
	\begin{bemerkung}
		Diese Vermutung markiert den entscheidenden Unterschied zum Zweifaktormodell. Dort fallen
		Balance und Eintritt oft auf den kleineren Faktor, während der Wechsel durch den größeren
		oder durch die gewichtete Log-Masse bestimmt wird. Im Dreifaktorfall scheint erstmals eine
		echte dreistufige Dynamik sichtbar zu werden.
	\end{bemerkung}
	
	\begin{hypothese}[Begrenzte Tragfähigkeit der Max-Log-Massen-Regel]
		\label{hyp:dreifaktor_maxlog_begrenzt}
		Im Zweifaktorfall \(p^a q^b\) bestimmt die größere gewichtete Log-Masse
		\[
		\max\{a\log p,\; b\log q\}
		\]
		numerisch zuverlässig die Wechselprimzahl. Im Dreifaktorfall
		\[
		p^a q^b r^c
		\]
		ist die entsprechende naive Verallgemeinerung
		\[
		\pi_{\mathrm{wech}}^{M}(n)=\arg\max\{a\log p,\; b\log q,\; c\log r\}
		\]
		im Allgemeinen nicht mehr hinreichend.
		
		Stattdessen ist numerisch plausibel, dass die Wechselprimzahl durch den \emph{letzten
			dominanten Restblock} bestimmt wird, der nach dem Eintritt früherer Faktoren in die
		Schale noch verbleibt.
	\end{hypothese}
	
	\begin{bemerkung}
		Dies würde erklären, warum die Max-Log-Massen-Regel im Dreifaktorfall zwar deutlich besser
		ist als die naive Regel \(\pi_{\mathrm{wech}}=r\), aber dennoch nicht vollständig
		zutrifft: Entscheidend ist offenbar nicht die größte Einzelmasse, sondern ihre Stellung
		innerhalb der fortschreitenden Schalenreduktion.
	\end{bemerkung}
	
	\begin{hypothese}[Erste Spezialisierung für den Basisfall \(pqr\)]
		\label{hyp:dreifaktor_pqr}
		Im ungewichteten Basisfall
		\[
		n=pqr,\qquad p<q<r,
		\]
		ist numerisch besonders plausibel:
		
		\begin{enumerate}
			\item Die Balanceprimzahl liegt häufig bei
			\[
			\pi_{\mathrm{bal}}^{M}(n)=q.
			\]
			\item Die Wechselprimzahl liegt häufig entweder bei \(q\) oder bei \(r\), abhängig davon,
			wie stark der letzte Restblock nach Eintritt von \(p\) und \(q\) noch dominiert.
		\end{enumerate}
	\end{hypothese}
	
	\begin{beispiel}
		Die bisherigen numerischen Beispiele stützen diese Sichtweise:
		
		\begin{enumerate}
			\item Für
			\[
			n=5\cdot 7\cdot 11
			\]
			wird numerisch beobachtet:
			\[
			\pi_{\mathrm{bal}}(n)=7,\qquad \pi_{\mathrm{wech}}(n)=7.
			\]
			
			\item Für
			\[
			n=5\cdot 7\cdot 11^2
			\]
			wird numerisch beobachtet:
			\[
			\pi_{\mathrm{bal}}(n)=7,\qquad \pi_{\mathrm{wech}}(n)=11.
			\]
			
			\item Für
			\[
			n=5^2\cdot 7\cdot 11
			\]
			wird numerisch beobachtet:
			\[
			\pi_{\mathrm{bal}}(n)=5,\qquad \pi_{\mathrm{wech}}(n)=7.
			\]
		\end{enumerate}
		
		Diese Beispiele sprechen dafür, dass Exponenten die Rollen der Faktoren gezielt
		verschieben können, ohne dass die Balance- und Wechselstruktur kollabiert.
	\end{beispiel}
	
	\begin{bemerkung}
		Die Ergebnisse deuten darauf hin, dass der Dreifaktorfall nicht mehr durch eine bloße
		Klein-Groß- oder Max-Log-Regel beschrieben werden kann. Vielmehr scheint hier erstmals
		eine genuine mehrstufige Schalen-Dynamik aufzutreten, in der mittlere Faktoren eine
		eigenständige strukturelle Rolle spielen.
	\end{bemerkung}
	
	\subsection{Schalenabschluss und Geburt neuer Zahlen}
	
	Wir formulieren nun präzise, wann eine gegebene Primschale die natürlichen Zahlen
	noch vollständig ausfüllen kann und ab welchem Punkt notwendig neue Primträger
	auftreten müssen.
	
	\begin{definition}[Primschale und Schalenmonoid]
		Für \(k\in\mathbb N\) sei
		\[
		P_k:=\{2,3,5,\dots,p_k\}
		\]
		die Menge der ersten \(k\) Primzahlen. Das von \(P_k\) erzeugte multiplikative Monoid sei
		\[
		\langle P_k\rangle
		:=
		\left\{
		\prod_{p\in P_k} p^{\nu_p}
		\;\middle|\;
		\nu_p\in\mathbb N_0
		\right\}.
		\]
		Die Elemente von \(\langle P_k\rangle\) heißen \(P_k\)-glatte Zahlen.
	\end{definition}
	
	\begin{definition}[Schalenanteil und reduzierter Kern]
		Für \(n\in\mathbb N\) definieren wir den Schalenanteil bezüglich \(P_k\) durch
		\[
		S_k(n):=\prod_{p\in P_k} p^{v_p(n)}
		\]
		und den reduzierten Kern durch
		\[
		R_k(n):=\frac{n}{S_k(n)}.
		\]
		Dabei enthält \(S_k(n)\) genau den aus den vorhandenen Schalenprimzahlen gebildeten Anteil
		von \(n\), während \(R_k(n)\) den außerhalb der Schale verbleibenden Rest beschreibt.
	\end{definition}
	
	\begin{lemma}[Charakterisierung des Schalenabschlusses]
		\label{lem:schalenabschluss}
		Für \(n\in\mathbb N\) gilt:
		\[
		n\in \langle P_k\rangle
		\quad\Longleftrightarrow\quad
		R_k(n)=1.
		\]
	\end{lemma}
	
	\begin{proof}
		Per Definition ist
		\[
		S_k(n)=\prod_{p\in P_k} p^{v_p(n)}
		\]
		genau der Anteil von \(n\), der aus Primfaktoren in \(P_k\) besteht.
		
		\medskip
		\noindent
		Falls \(n\in \langle P_k\rangle\), so besitzt \(n\) keine Primfaktoren außerhalb von \(P_k\).
		Dann enthält \(S_k(n)\) bereits die gesamte Primfaktorzerlegung von \(n\), also
		\[
		S_k(n)=n
		\qquad\Longrightarrow\qquad
		R_k(n)=\frac{n}{S_k(n)}=1.
		\]
		
		\medskip
		\noindent
		Umgekehrt, falls \(R_k(n)=1\), gilt
		\[
		n=S_k(n),
		\]
		also besteht \(n\) vollständig aus Primfaktoren in \(P_k\). Somit ist
		\[
		n\in \langle P_k\rangle.
		\]
	\end{proof}
	
	\begin{satz}[Scharfe Grenze der lückenlosen Ausfüllung]
		\label{satz:scharfe_grenze_schale}
		Für die Primschale
		\[
		P_k=\{2,3,\dots,p_k\}
		\]
		gilt:
		\[
		\{1,2,\dots,p_{k+1}-1\}\subseteq \langle P_k\rangle,
		\]
		während
		\[
		p_{k+1}\notin \langle P_k\rangle.
		\]
		
		Insbesondere ist \(p_{k+1}\) die erste natürliche Zahl, die mit den bereits
		vorhandenen Primfaktoren aus \(P_k\) nicht mehr dargestellt werden kann.
	\end{satz}
	
	\begin{proof}
		Sei \(n<p_{k+1}\). Dann kann kein Primfaktor von \(n\) größer oder gleich \(p_{k+1}\) sein,
		denn sonst wäre dieser Primfaktor bereits größer als \(n\), was unmöglich ist. Also haben
		alle Primfaktoren von \(n\) die Form
		\[
		p\le p_k.
		\]
		Damit besteht die Primfaktorzerlegung von \(n\) ausschließlich aus Primzahlen in \(P_k\), also
		\[
		n\in \langle P_k\rangle.
		\]
		
		Andererseits ist \(p_{k+1}\) selbst prim und gehört nicht zu \(P_k\). Folglich kann
		\(p_{k+1}\) nicht als Produkt von Primzahlen aus \(P_k\) dargestellt werden, also
		\[
		p_{k+1}\notin \langle P_k\rangle.
		\]
		Damit ist \(p_{k+1}\) die erste Lücke.
	\end{proof}
	
	\begin{korollar}[Geburt eines neuen Primträgers]
		\label{kor:geburt_neuer_primtraeger}
		Für eine gegebene Zahl \(n\in\mathbb N\) relativ zur Schale \(P_k\) gilt:
		\[
		R_k(n)>1
		\quad\Longleftrightarrow\quad
		\text{in } n \text{ tritt mindestens ein neuer Primfaktor } >p_k \text{ auf.}
		\]
	\end{korollar}
	
	\begin{proof}
		Aus Lemma~\ref{lem:schalenabschluss} folgt
		\[
		R_k(n)=1
		\quad\Longleftrightarrow\quad
		n\in\langle P_k\rangle.
		\]
		Also gilt äquivalent:
		\[
		R_k(n)>1
		\quad\Longleftrightarrow\quad
		n\notin \langle P_k\rangle.
		\]
		Dies bedeutet genau, dass \(n\) mindestens einen Primfaktor besitzt, der nicht in \(P_k\)
		liegt, also größer als \(p_k\) ist.
	\end{proof}
	
	\begin{definition}[E/ABC-Zerlegung des reduzierten Kerns]
		Wir zerlegen den reduzierten Kern nach den Restklassen modulo \(12\):
		\[
		R_k(n)=E_k(n)\,A_k(n)\,B_k(n)\,C_k(n),
		\]
		wobei
		\[
		E_k(n):=\prod_{\substack{q^a\parallel R_k(n)\\ q\equiv 1\pmod{12}}} q^a,
		\]
		\[
		A_k(n):=\prod_{\substack{q^a\parallel R_k(n)\\ q\equiv 5\pmod{12}}} q^a,\qquad
		B_k(n):=\prod_{\substack{q^a\parallel R_k(n)\\ q\equiv 7\pmod{12}}} q^a,\qquad
		C_k(n):=\prod_{\substack{q^a\parallel R_k(n)\\ q\equiv 11\pmod{12}}} q^a.
		\]
	\end{definition}
	
	\begin{satz}[Kriterium für das Auftreten neuer Zahlen in E/ABC-Sprache]
		\label{satz:eabc_kriterium_neue_zahl}
		Für \(n\in\mathbb N\) relativ zur Schale \(P_k\) gilt:
		\[
		R_k(n)>1
		\quad\Longleftrightarrow\quad
		E_k(n)A_k(n)B_k(n)C_k(n)>1.
		\]
		
		Damit liefert die E/ABC-Zerlegung des reduzierten Kerns kein neues Existenzkriterium,
		wohl aber die \emph{Signatur} des neuen Primträgers:
		\begin{itemize}
			\item \(E_k(n)>1\), \(A_k(n)=B_k(n)=C_k(n)=1\): rein energetischer neuer Rest,
			\item \(E_k(n)=1\), \(A_k(n)B_k(n)C_k(n)>1\): rein \(ABC\)-artiger neuer Rest,
			\item \(E_k(n)>1\) und \(A_k(n)B_k(n)C_k(n)>1\): gemischter neuer Rest.
		\end{itemize}
	\end{satz}
	
	\begin{proof}
		Da per Definition
		\[
		R_k(n)=E_k(n)\,A_k(n)\,B_k(n)\,C_k(n),
		\]
		gilt unmittelbar:
		\[
		R_k(n)>1
		\quad\Longleftrightarrow\quad
		E_k(n)A_k(n)B_k(n)C_k(n)>1.
		\]
		Die Fallunterscheidung beschreibt lediglich, in welchen Restklassen die neuen
		Primfaktoren des Kerns auftreten.
	\end{proof}
	
	\begin{bemerkung}
		Die Grenze
		\[
		p_{k+1}-1
		\]
		ist die schärfste mögliche Grenze der \emph{lückenlosen} Ausfüllung der natürlichen Zahlen
		durch die vorhandene Schale. Ab \(p_{k+1}\) bricht nicht die Existenz weiterer
		\(P_k\)-glatter Zahlen zusammen, wohl aber ihre Vollständigkeit als Anfangsstrecke.
	\end{bemerkung}
	
	\begin{bemerkung}
		In der Sprache des Schalenwechsel-Modells bedeutet dies: Die Schale ist genau dann
		abgeschlossen, wenn der reduzierte Kern trivial ist. Die Geburt einer neuen Zahl im
		arithmetisch relevanten Sinn beginnt genau dann, wenn
		\[
		R_k(n)>1.
		\]
		Die E/ABC-Zerlegung entscheidet anschließend nicht mehr über das bloße Auftreten des
		Neuen, sondern über dessen innere Signatur.
	\end{bemerkung}
	
	\subsection{Numerischer Test bis \(10^7\)}
	
	Zur numerischen Überprüfung der Aussagen über Schalenabschluss und reduzierte Kerne
	wurde ein vollständiger Test für alle natürlichen Zahlen
	\[
	1\le n\le 10^7
	\]
	durchgeführt. Grundlage war ein Siebverfahren, das für jede Zahl \(n\) den größten
	Primfaktor
	\[
	P^+(n)
	\]
	bestimmt. Damit ist die Zugehörigkeit zum von einer Schale \(P_k=\{2,3,\dots,p_k\}\)
	erzeugten Monoid äquivalent zu
	\[
	n\in \langle P_k\rangle
	\quad\Longleftrightarrow\quad
	P^+(n)\le p_k.
	\]
	Nach Lemma~\ref{lem:schalenabschluss} ist dies wiederum äquivalent zu
	\[
	R_k(n)=1.
	\]
	
	\begin{satz}[Numerische Bestätigung bis \(10^7\)]
		\label{satz:numerik_bis_10hoch7}
		Für alle getesteten Schalen \(P_k\) mit
		\[
		p_{k+1}\le 10^7
		\]
		und für alle Zahlen
		\[
		1\le n\le 10^7
		\]
		bestätigen die numerischen Daten:
		
		\begin{enumerate}
			\item Die Anfangsstrecke
			\[
			\{1,2,\dots,p_{k+1}-1\}
			\]
			liegt vollständig in \(\langle P_k\rangle\).
			
			\item Die erste nicht mehr darstellbare Zahl ist genau
			\[
			p_{k+1}.
			\]
			
			\item Für jede getestete Zahl \(n\) gilt äquivalent:
			\[
			R_k(n)=1
			\quad\Longleftrightarrow\quad
			P^+(n)\le p_k,
			\]
			also
			\[
			R_k(n)>1
			\quad\Longleftrightarrow\quad
			P^+(n)>p_k.
			\]
		\end{enumerate}
	\end{satz}
	
	\begin{proof}[Numerische Verifikation]
		Der Test wurde vollständig für alle \(n\le 10^7\) ausgeführt. Für jede Zahl wurde der
		größte Primfaktor \(P^+(n)\) berechnet. Damit ließ sich für jede Schale \(P_k\) direkt
		prüfen, ob \(n\) vollständig aus den vorhandenen Schalenprimzahlen aufgebaut werden kann.
		Dabei traten keine Gegenbeispiele zu Satz~\ref{satz:scharfe_grenze_schale},
		Lemma~\ref{lem:schalenabschluss} oder Korollar~\ref{kor:geburt_neuer_primtraeger} auf.
	\end{proof}
	
	\begin{beispiel}
		Die Anzahl der \(P_k\)-glatten Zahlen bis \(10^7\) beträgt für einige ausgewählte Schalen:
		
		\[
		\#\{n\le 10^7:\ P^+(n)\le 2\}=24,
		\]
		\[
		\#\{n\le 10^7:\ P^+(n)\le 3\}=190,
		\]
		\[
		\#\{n\le 10^7:\ P^+(n)\le 5\}=768,
		\]
		\[
		\#\{n\le 10^7:\ P^+(n)\le 7\}=2155,
		\]
		\[
		\#\{n\le 10^7:\ P^+(n)\le 11\}=4520,
		\]
		\[
		\#\{n\le 10^7:\ P^+(n)\le 13\}=8289,
		\]
		\[
		\#\{n\le 10^7:\ P^+(n)\le 29\}=37627,
		\]
		\[
		\#\{n\le 10^7:\ P^+(n)\le 97\}=269882.
		\]
		
		Diese Zahlen zeigen, dass auch für größere Schalen noch viele glatte Zahlen auftreten,
		dass ihre Dichte aber rasch abnimmt.
	\end{beispiel}
	
	\begin{bemerkung}
		Die numerischen Tests bestätigen damit sehr deutlich die begriffliche Trennung zwischen
		\emph{Schalenabschluss} und \emph{Signatur des Neuen}. Das Existenzkriterium für einen
		neuen Primträger ist vollständig durch
		\[
		R_k(n)>1
		\]
		gegeben. Die E/ABC-Zerlegung von \(R_k(n)\) entscheidet anschließend nicht mehr darüber,
		\emph{ob} etwas Neues auftritt, sondern nur darüber, \emph{welcher Art} dieses Neue ist.
	\end{bemerkung}
	
	\begin{bemerkung}
		Die Grenze
		\[
		p_{k+1}-1
		\]
		ist numerisch bis \(10^7\) ausnahmslos als scharfe Grenze der lückenlosen Ausfüllung
		bestätigt worden. Ab \(p_{k+1}\) bricht nicht die Existenz weiterer \(P_k\)-glatter Zahlen
		zusammen, wohl aber ihre Vollständigkeit als Anfangsstrecke.
	\end{bemerkung}
	
	\subsection{Schalenabschluss und Primzahlvierlinge}
	
	Die Aussagen über Schalenabschluss und reduzierte Kerne besitzen eine natürliche lokale
	Ergänzung durch die Struktur der Primzahlvierlinge. Während der Schalenabschluss global
	entscheidet, bis zu welchem Punkt die natürlichen Zahlen durch eine gegebene Schale
	vollständig erzeugt werden können, beschreiben Primzahlvierlinge die dichteste lokale
	Organisation neuer Primstruktur jenseits dieser Schale.
	
	\begin{definition}[Primzahlvierling]
		Ein Primzahlvierling ist eine Primzahlkonstellation der Form
		\[
		(p,p+2,p+6,p+8),
		\]
		bei der alle vier Zahlen prim sind.
	\end{definition}
	
	\begin{bemerkung}
		Für Primzahlen \(>3\) sind modulo \(12\) nur die Restklassen
		\[
		1,5,7,11
		\]
		möglich. In der Sprache der Familien
		\[
		E,\ A,\ B,\ C
		\]
		entspricht dies genau den vier zugelassenen Signaturen jenseits der \(2\)- und
		\(3\)-Schale.
	\end{bemerkung}
	
	\begin{lemma}[EABC-Zyklus im Primzahlvierling]
		\label{lem:vierling_eabc}
		Sei
		\[
		(p,p+2,p+6,p+8)
		\]
		ein Primzahlvierling mit \(p>3\). Dann durchlaufen die vier Einträge modulo \(12\) genau
		einmal die vier Restklassen
		\[
		5,\ 7,\ 11,\ 1,
		\]
		also in Familiennotation
		\[
		A,\ B,\ C,\ E,
		\]
		gegebenenfalls zyklisch verschoben.
	\end{lemma}
	
	\begin{proof}
		Da \(p>3\) prim ist, kann \(p\) modulo \(12\) nur in einer der Klassen
		\[
		1,\ 5,\ 7,\ 11
		\]
		liegen. Addiert man die Verschiebungen \(0,2,6,8\), so erhält man modulo \(12\) genau die
		vier Werte
		\[
		p,\ p+2,\ p+6,\ p+8
		\]
		als vollständige zyklische Durchläufe durch die vier zulässigen Primrestklassen. Für den
		Startwert \(p\equiv 5\pmod{12}\) ergibt sich unmittelbar
		\[
		5,\ 7,\ 11,\ 1.
		\]
		Die übrigen Fälle entstehen durch zyklische Verschiebung derselben Viererstruktur.
	\end{proof}
	
	\begin{satz}[Globale und lokale Seite neuer Primstruktur]
		\label{satz:global_lokal_neu}
		Sei \(P_k=\{2,3,\dots,p_k\}\) eine Primschale. Dann gilt:
		
		\begin{enumerate}
			\item Global wird das Auftreten neuer Primstruktur relativ zu \(P_k\) exakt durch den
			reduzierten Kern entschieden:
			\[
			R_k(n)>1
			\quad\Longleftrightarrow\quad
			\text{in } n \text{ tritt mindestens ein Primfaktor } >p_k \text{ auf.}
			\]
			
			\item Lokal liefern Primzahlvierlinge die dichteste Konfiguration, in der neue
			Primstruktur jenseits der \(2\)- und \(3\)-Schale als vollständiger \(EABC\)-Block
			sichtbar wird.
		\end{enumerate}
	\end{satz}
	
	\begin{proof}[Begründung]
		Der erste Teil ist bereits in Korollar~\ref{kor:geburt_neuer_primtraeger} gezeigt.
		
		Für den zweiten Teil ist entscheidend, dass Primzahlen \(>3\) nur in den vier Klassen
		\[
		1,\ 5,\ 7,\ 11\pmod{12}
		\]
		auftreten können. Ein Primzahlvierling realisiert diese vier Klassen in maximal dichter
		Anordnung, da die Form
		\[
		(p,p+2,p+6,p+8)
		\]
		die engste mögliche Konfiguration ist, die mit den Restriktionen der Teilbarkeit durch
		\(2\) und \(3\) verträglich bleibt. Nach Lemma~\ref{lem:vierling_eabc} erscheint dabei
		lokal genau einmal die vollständige EABC-Signatur.
	\end{proof}
	
	\begin{bemerkung}
		Damit ergibt sich eine konzeptionelle Trennung:
		
		\begin{itemize}
			\item Der reduzierte Kern \(R_k(n)\) beantwortet die globale Existenzfrage:
			\[
			\text{„Ist neue Primstruktur außerhalb der Schale vorhanden?“}
			\]
			
			\item Der Primzahlvierling beantwortet eine lokale Organisationsfrage:
			\[
			\text{„Wie kann neue Primstruktur in maximal dichter EABC-Form auftreten?“}
			\]
		\end{itemize}
		
		In diesem Sinn misst \(R_k(n)>1\) das Auftreten des Neuen, während der Primzahlvierling
		seine dichteste lokale EABC-Kohärenz beschreibt.
	\end{bemerkung}
	
	\begin{korollar}[Primzahlvierlinge als lokale Vollsignatur]
		\label{kor:vierling_vollsignatur}
		Primzahlvierlinge sind die kleinste lokale Primkonfiguration, in der alle vier Familien
		\[
		E,\ A,\ B,\ C
		\]
		gleichzeitig auftreten. Sie bilden damit eine lokale Vollsignatur neuer Primstruktur
		jenseits der \(2\)- und \(3\)-Schale.
	\end{korollar}
	
	\begin{bemerkung}
		Die Aussage über die erste nicht mehr ausfüllbare Zahl
		\[
		p_{k+1}
		\]
		beschreibt die erste globale Lücke einer gegebenen Schale. Primzahlvierlinge beschreiben
		demgegenüber nicht bloß das Auftreten einer einzelnen neuen Primzahl, sondern eine
		strukturierte Überschreitung, in der neue Primzahlen in maximal dichter und zugleich
		vollständiger EABC-Organisation erscheinen.
	\end{bemerkung}
	
	\begin{bemerkung}
		In der Sprache des Schalenwechsel-Modells kann man daher sagen: Der Schalenabschluss
		kontrolliert die globale Vollständigkeit der aus vorhandenen Primträgern erzeugbaren
		Anfangsstrecke, während Primzahlvierlinge die erste lokal maximale Gestalt des Neuen
		darstellen, sobald die Restriktionen der kleinen Schalen überwunden sind.
	\end{bemerkung}
	
	\subsection{Leitformel für Schalenabschluss und lokale Vierlingsstruktur}
	
	Die grundlegende Größe des Modells ist der reduzierte Kern relativ zu einer gegebenen
	Primschale
	\[
	P_k=\{2,3,5,\dots,p_k\}.
	\]
	Er wird durch
	\[
	\boxed{
		R_k(n):=\frac{n}{\displaystyle\prod_{p\in P_k} p^{v_p(n)}}
	}
	\]
	definiert.
	
	\begin{bemerkung}[Leitformel des Schalenmodells]
		\label{prop:leitformel_schale}
		Für jede natürliche Zahl \(n\) und jede Primschale \(P_k\) gilt:
		
		\[
		\boxed{
			R_k(n)=1
			\iff
			n\in \langle P_k\rangle
		}
		\]
		
		und äquivalent
		
		\[
		\boxed{
			R_k(n)>1
			\iff
			\text{in }n\text{ tritt mindestens ein Primfaktor }>p_k\text{ auf.}
		}
		\]
	\end{bemerkung}
	
	\begin{proof}
		Die Zahl
		\[
		S_k(n):=\prod_{p\in P_k} p^{v_p(n)}
		\]
		enthält per Definition genau alle Primfaktorpotenzen von \(n\), deren Primfaktoren bereits
		in der Schale \(P_k\) liegen. Daher ist
		\[
		R_k(n)=\frac{n}{S_k(n)}
		\]
		genau der außerhalb der Schale verbleibende Rest.
		
		\medskip
		\noindent
		Falls \(R_k(n)=1\), gilt \(n=S_k(n)\), also besteht \(n\) vollständig aus Primfaktoren in
		\(P_k\). Daher ist
		\[
		n\in \langle P_k\rangle.
		\]
		
		Umgekehrt, falls \(n\in \langle P_k\rangle\), besitzt \(n\) keine Primfaktoren außerhalb
		von \(P_k\), also ist bereits \(S_k(n)=n\), und damit
		\[
		R_k(n)=1.
		\]
		
		Die zweite Äquivalenz ist nur die Umformulierung der Negation: \(R_k(n)>1\) bedeutet
		genau, dass nach Abzug des Schalenanteils noch ein nichttrivialer Rest außerhalb der
		Schale verbleibt.
	\end{proof}
	
	\begin{korollar}[Scharfe Grenze der lückenlosen Ausfüllung]
		\label{kor:leitformel_grenze}
		Für die Primschale \(P_k\) ist die erste nicht mehr vollständig durch die vorhandenen
		Primfaktoren darstellbare Zahl genau
		\[
		p_{k+1}.
		\]
		Insbesondere gilt
		\[
		\{1,2,\dots,p_{k+1}-1\}\subseteq \langle P_k\rangle,
		\qquad
		p_{k+1}\notin \langle P_k\rangle.
		\]
	\end{korollar}
	
	\begin{proof}
		Jede Zahl \(n<p_{k+1}\) kann keinen Primfaktor \(\ge p_{k+1}\) besitzen und liegt daher
		in \(\langle P_k\rangle\). Die Zahl \(p_{k+1}\) selbst ist prim und gehört nicht zu
		\(P_k\), also nicht zu \(\langle P_k\rangle\).
	\end{proof}
	
	\begin{definition}[Vierlingsindikator]
		Für \(p\in\mathbb N\) definieren wir den Indikator eines Primzahlvierlings durch
		\[
		\boxed{
			\mathcal V(p):=
			\mathbf 1_{\mathbb P}(p)\,
			\mathbf 1_{\mathbb P}(p+2)\,
			\mathbf 1_{\mathbb P}(p+6)\,
			\mathbf 1_{\mathbb P}(p+8)
		}
		\]
		wobei \(\mathbf 1_{\mathbb P}\) die Indikatorfunktion der Primzahlen bezeichnet.
	\end{definition}
	
	\begin{korollar}[Lokale Vollsignatur eines Vierlings]
		\label{kor:leitformel_vierling}
		Falls
		\[
		\mathcal V(p)=1,
		\]
		also
		\[
		(p,p+2,p+6,p+8)
		\]
		ein Primzahlvierling ist, dann durchlaufen diese vier Zahlen modulo \(12\) genau die
		vier Primrestklassen
		\[
		5,\ 7,\ 11,\ 1,
		\]
		also in Familiennotation
		\[
		A,\ B,\ C,\ E,
		\]
		bis auf zyklische Verschiebung.
	\end{korollar}
	
	\begin{proof}
		Für Primzahlen \(>3\) sind modulo \(12\) nur die Klassen
		\[
		1,\ 5,\ 7,\ 11
		\]
		möglich. Die Verschiebungen \(0,2,6,8\) erzeugen modulo \(12\) genau einen vollständigen
		zyklischen Durchlauf durch diese vier Klassen.
	\end{proof}
	
	\begin{bemerkung}
		Die Leitformel
		\[
		R_k(n)=\frac{n}{S_k(n)}
		\]
		misst global, ob eine Zahl bereits vollständig durch die Schale erzeugt werden kann.
		Der Vierlingsindikator
		\[
		\mathcal V(p)
		\]
		misst dagegen lokal eine maximal dichte Primkonfiguration. In diesem Sinn gilt:
		
		\[
		\boxed{
			R_k(n)>1
			\text{ misst die Existenz des Neuen,}
			\qquad
			\mathcal V(p)=1
			\text{ misst seine dichteste lokale EABC-Ausprägung.}
		}
		\]
	\end{bemerkung}
	
	\subsection{Die EA/BC-Zweischale als kleinste raue Schalenstruktur}
	
	Nach Abtrennung der glatten Grundkoeffizienten \(2\), \(3\) und \(5\) bleibt ein rauer
	Kern zurück, auf dem sich die vier Familien \(E,A,B,C\) zu einer kleinsten nichttrivialen
	Zweischale verdichten. Diese Zweischale wird durch die Paarung
	\[
	EA
	\qquad\text{und}\qquad
	BC
	\]
	gegeben.
	
	\begin{definition}[Rauer Kern]
		Für \(n\in\mathbb N\) definieren wir den von den glatten Faktoren \(2\), \(3\) und \(5\)
		befreiten rauen Kern durch
		\[
		m(n):=\frac{n}{2^{v_2(n)}3^{v_3(n)}5^{v_5(n)}}.
		\]
		Dann gilt stets
		\[
		\gcd(m(n),30)=1.
		\]
		Insbesondere liegen alle Primfaktoren von \(m(n)\) in den Restklassen
		\[
		1,\ 5,\ 7,\ 11 \pmod{12},
		\]
		also in den Familien \(E,A,B,C\).
	\end{definition}
	
	\begin{definition}[EA/BC-Signatur]
		\label{def:eabc_zweischale}
		Für einen rauen Kern \(m\) mit \(\gcd(m,30)=1\) definieren wir die Zweischalen-Signatur
		\[
		\chi_{EA/BC}(m):=
		\begin{cases}
			+1,& m\bmod 12\in\{1,5\},\\[1mm]
			-1,& m\bmod 12\in\{7,11\}.
		\end{cases}
		\]
		Wir sagen:
		\[
		m\in EA
		\iff
		\chi_{EA/BC}(m)=+1,
		\]
		\[
		m\in BC
		\iff
		\chi_{EA/BC}(m)=-1.
		\]
	\end{definition}
	
	\begin{lemma}[EA ist Untergruppe, BC ist Nebenklasse]
		\label{lem:eabc_untergruppe}
		Die Menge
		\[
		\{1,5\}\subset (\mathbb Z/12\mathbb Z)^\times
		\]
		ist eine Untergruppe der Einheitengruppe modulo \(12\), und
		\[
		\{7,11\}
		\]
		ist die komplementäre Nebenklasse. Insbesondere gilt auf Signaturniveau
		\[
		EA\cdot EA=EA,\qquad
		EA\cdot BC=BC,\qquad
		BC\cdot BC=EA.
		\]
	\end{lemma}
	
	\begin{proof}
		Die Einheiten modulo \(12\) sind
		\[
		(\mathbb Z/12\mathbb Z)^\times=\{1,5,7,11\}.
		\]
		Man rechnet unmittelbar nach:
		\[
		1\cdot 1\equiv 1,\qquad
		1\cdot 5\equiv 5,\qquad
		5\cdot 5\equiv 1 \pmod{12}.
		\]
		Damit ist \(\{1,5\}\) eine Untergruppe.
		
		Ebenso gilt
		\[
		7\cdot 7\equiv 1,\qquad
		7\cdot 11\equiv 5,\qquad
		11\cdot 11\equiv 1 \pmod{12},
		\]
		sodass das Produkt zweier Elemente aus \(\{7,11\}\) wieder in \(\{1,5\}\) liegt. Weiter
		gilt
		\[
		1\cdot 7\equiv 7,\qquad
		1\cdot 11\equiv 11,\qquad
		5\cdot 7\equiv 11,\qquad
		5\cdot 11\equiv 7 \pmod{12}.
		\]
		Damit folgt genau die angegebene Zweischalen-Algebra.
	\end{proof}
	
	\begin{satz}[Faktorisierungskriterium der EA/BC-Zweischale]
		\label{satz:eabc_faktorisierung}
		Sei
		\[
		m=\prod_{i=1}^r p_i^{e_i}
		\]
		die Primfaktorzerlegung eines rauen Kerns \(m\), also \(\gcd(m,30)=1\). Dann gilt
		\[
		\boxed{
			\chi_{EA/BC}(m)
			=
			(-1)^{\sum_{p_i\in B\cup C} e_i}.
		}
		\]
		
		Insbesondere gilt:
		\begin{itemize}
			\item Eine gerade Gesamtzahl von Exponenten aus den Familien \(B\) und \(C\) impliziert
			\[
			m\in EA.
			\]
			\item Eine ungerade Gesamtzahl von Exponenten aus den Familien \(B\) und \(C\) impliziert
			\[
			m\in BC.
			\]
		\end{itemize}
	\end{satz}
	
	\begin{proof}
		Auf Zweischalen-Niveau wirken Faktoren aus \(E\) und \(A\) neutral, also mit Vorzeichen
		\(+1\), während Faktoren aus \(B\) und \(C\) jeweils einen Vorzeichenwechsel erzeugen.
		Da die Multiplikation in der Zweischale nach Lemma~\ref{lem:eabc_untergruppe} genau der
		Multiplikation von \(\pm1\) entspricht, liefert jede Primfaktorpotenz \(p_i^{e_i}\) mit
		\(p_i\in B\cup C\) den Beitrag \((-1)^{e_i}\). Multiplikation aller Beiträge ergibt
		\[
		\chi_{EA/BC}(m)
		=
		\prod_{p_i\in B\cup C} (-1)^{e_i}
		=
		(-1)^{\sum_{p_i\in B\cup C} e_i}.
		\]
	\end{proof}
	
	\begin{korollar}[Kleinste nichttriviale Schalenstruktur]
		\label{kor:eabc_kleinste_schale}
		Nach Abtrennung der glatten Faktoren \(2\), \(3\) und \(5\) ist die Paarung
		\[
		EA\mid BC
		\]
		die kleinste nichttriviale Schalenstruktur auf dem rauen Kern. Sie ist zweistufig und wird
		vollständig durch die Signatur
		\[
		\chi_{EA/BC}(m)\in\{+1,-1\}
		\]
		beschrieben.
	\end{korollar}
	
	\begin{bemerkung}
		Die feine Viererstruktur \(E,A,B,C\) wird in der EA/BC-Zweischale auf eine gröbere, aber
		algebraisch sehr stabile Zwei-Klassen-Struktur projiziert. Diese Projektion bewahrt den
		wesentlichen Flip-Mechanismus:
		\[
		B\ \text{oder}\ C \text{ wechseln die Zweischale,}
		\qquad
		E\ \text{oder}\ A \text{lassen sie unverändert.}
		\]
	\end{bemerkung}
	
	\begin{definition}[Verfeinerte raue Schalenkoordinate]
		Für einen rauen Kern \(m\) definieren wir die kombinierte Schalenkoordinate
		\[
		\boxed{
			\bigl(P^+(m),\chi_{EA/BC}(m)\bigr),
		}
		\]
		wobei \(P^+(m)\) den größten Primfaktor von \(m\) bezeichnet.
	\end{definition}
	
	\begin{bemerkung}
		Die Größe \(P^+(m)\) misst, in welcher rauen Größenschale der Kern liegt, während
		\[
		\chi_{EA/BC}(m)
		\]
		angibt, in welchem der beiden kleinsten Paarsektoren \(EA\) bzw. \(BC\) er sich befindet.
		Damit erhält man eine zweidimensionale Schalenbeschreibung:
		\[
		\text{Größenlage}\times\text{Paar-Signatur}.
		\]
	\end{bemerkung}
	
	\begin{korollar}[Primzahlvierlinge in der EA/BC-Projektion]
		\label{kor:vierling_eabc_zweischale}
		Ein Primzahlvierling
		\[
		(p,p+2,p+6,p+8)
		\]
		mit \(p>3\) besitzt in der feinen Familienstruktur die Signatur
		\[
		(A,B,C,E)
		\]
		bis auf zyklische Verschiebung. In der groben EA/BC-Zweischale projiziert sich dies zu
		\[
		(EA,\,BC,\,BC,\,EA)
		\]
		bis auf zyklische Verschiebung.
	\end{korollar}
	
	\begin{proof}
		Die feine Signatur folgt aus der modulo-\(12\)-Struktur des Vierlings. Unter der Paarung
		\[
		E,A\mapsto EA,
		\qquad
		B,C\mapsto BC
		\]
		wird daraus unmittelbar
		\[
		(A,B,C,E)\mapsto (EA,BC,BC,EA).
		\]
	\end{proof}
	
	\begin{bemerkung}
		Damit erscheinen Primzahlvierlinge in der groben Zweischalenprojektion als symmetrische
		Zwei-Sektoren-Muster mit einem \(BC\)-Kern und \(EA\)-Rand. Dies spricht dafür, dass die
		EA/BC-Zweischale tatsächlich die kleinste sinnvolle Projektionsschale ist, in der
		lokale Primstruktur noch nicht trivial wird.
	\end{bemerkung}
	
	\subsection{Numerische Kopplung der Vierer- und Zweierschale}
	
	Zur Untersuchung der Beziehung zwischen der feinen Viererschale
	\[
	E,\ A,\ B,\ C
	\]
	und der groben Zweierschale
	\[
	EA\mid BC
	\]
	wurde für alle natürlichen Zahlen
	\[
	1\le n\le 10^6
	\]
	der raue Kern
	\[
	m(n):=\frac{n}{2^{v_2(n)}3^{v_3(n)}5^{v_5(n)}}
	\]
	berechnet und anschließend sowohl fein als auch grob klassifiziert.
	
	Die feine Klassifikation erfolgt über
	\[
	m(n)\bmod 12\in\{1,5,7,11\},
	\]
	also
	\[
	1\mapsto E,\qquad
	5\mapsto A,\qquad
	7\mapsto B,\qquad
	11\mapsto C.
	\]
	Die grobe Klassifikation ist die Projektion
	\[
	E,A\mapsto EA,
	\qquad
	B,C\mapsto BC.
	\]
	
	\begin{satz}[Numerische Kopplung bis \(10^6\)]
		\label{satz:numerik_kopplung_vierer_zweier}
		Für die rauen Kerne \(m(n)\) bis \(10^6\) zeigen die numerischen Daten:
		
		\begin{enumerate}
			\item Die feinen Klassen \(E,A,B,C\) sind praktisch gleichverteilt:
			\[
			P(E)\approx P(A)\approx P(B)\approx P(C)\approx \frac14.
			\]
			
			\item Die groben Klassen \(EA,BC\) sind praktisch gleichverteilt:
			\[
			P(EA)\approx P(BC)\approx \frac12.
			\]
			
			\item Die Shannon-Entropien erfüllen numerisch
			\[
			H_4\approx \log 4,
			\qquad
			H_2\approx \log 2,
			\qquad
			H_4-H_2\approx \log 2.
			\]
			
			\item Alle bis \(10^6\) gefundenen Primzahlvierlinge projizieren exakt auf das
			erwartete Zweischalenmuster
			\[
			(EA,BC,BC,EA)
			\]
			bis auf zyklische Verschiebung.
		\end{enumerate}
	\end{satz}
	
	\begin{proof}[Numerische Verifikation]
		Die numerische Auswertung ergab für die feinen Klassenverteilungen:
		
		\[
		E: 249753,\qquad
		A: 249813,\qquad
		B: 250077,\qquad
		C: 249850,
		\]
		also insgesamt
		\[
		999493
		\]
		nichttriviale raue Kerne. Für die groben Klassen ergab sich:
		
		\[
		EA: 499566,\qquad
		BC: 499927.
		\]
		
		Die zugehörigen Entropien lauten numerisch:
		\[
		H_4=1.386294240726,
		\qquad
		H_2=0.693147115333,
		\]
		also
		\[
		H_4-H_2=0.693147125392.
		\]
		Dies stimmt numerisch mit
		\[
		\log 4,\qquad \log 2,\qquad \log 2
		\]
		überein.
		
		Zusätzlich wurden bis \(10^6\) genau
		\[
		166
		\]
		Primzahlvierlinge gefunden. In allen \(166\) Fällen stimmte die Projektion auf die
		Zweierschale exakt mit dem Muster
		\[
		(EA,BC,BC,EA)
		\]
		bis auf zyklische Verschiebung überein.
	\end{proof}
	
	\begin{bemerkung}
		Die Zahlen
		\[
		H_4\approx \log 4,
		\qquad
		H_2\approx \log 2
		\]
		zeigen, dass sowohl die Viererschale als auch ihre Zweierschalenprojektion global nahezu
		maximal gleichverteilt sind. Damit erklärt sich die Zweierschale nicht als grobe
		Verzerrung, sondern als symmetrischer Quotient der feinen Viererstruktur.
	\end{bemerkung}
	
	\begin{bemerkung}
		Der Entropieverlust
		\[
		\Delta H:=H_4-H_2\approx \log 2
		\]
		hat eine klare informations-theoretische Bedeutung: Beim Übergang von
		\[
		E,A,B,C
		\]
		zu
		\[
		EA,BC
		\]
		geht genau eine binäre Entscheidungsstufe verloren, nämlich die innere Unterscheidung
		\[
		E\leftrightarrow A
		\qquad\text{und}\qquad
		B\leftrightarrow C.
		\]
		Die globale Symmetrie bleibt dabei erhalten.
	\end{bemerkung}
	
	\begin{bemerkung}
		Auch die Übergangsmatrizen bestätigen diese Interpretation. Auf grober Ebene verhält sich
		die Folge der Zweischalen-Signaturen nahezu wie ein symmetrischer Zwei-Zustands-Prozess,
		während die feine Viererschale eine deutlich reichere interne Übergangsstruktur trägt.
		Damit erscheint die Zweierschale als stabile Quotientenschale, die den globalen Sektor
		bewahrt, während die Viererschale die innere Orientierung auflöst.
	\end{bemerkung}
	
	\begin{korollar}[Kanonische Quotientenschale]
		\label{kor:kanonische_quotientenschale}
		Die EA/BC-Zweierschale ist numerisch die kanonische grobe Projektion der EABC-Viererschale.
		Sie bewahrt die globale Symmetrie vollständig, verliert dabei aber genau ein Bit an
		feiner Orientierungsinformation.
	\end{korollar}
	
	\begin{bemerkung}
		Die exakte Projektion aller getesteten Primzahlvierlinge auf das Zweischalenmuster
		\[
		(EA,BC,BC,EA)
		\]
		liefert darüber hinaus einen starken Hinweis darauf, dass die Zweierschale nicht bloß
		statistisch sinnvoll, sondern auch lokal mit der feineren Primkonstellationsgeometrie
		verträglich ist.
	\end{bemerkung}
	
	\subsection{Glatte Schalen und raue Signaturen}
	
	Die bisher betrachteten Schalenstrukturen lassen sich besonders klar verstehen, wenn man
	zwischen einem glatten Träger und einer rauen Signatur unterscheidet. Die glatten Zahlen
	bilden die radiale Schalenlage, während der nach Abzug der glatten Faktoren verbleibende
	raue Kern die eigentliche EABC- bzw. EA/BC-Orientierung trägt.
	
	\begin{definition}[Glatte Schale relativ zu \(2,3,5\)]
		Für \(n\in\mathbb N\) definieren wir den \((2,3,5)\)-glatten Anteil durch
		\[
		s_{235}(n):=2^{v_2(n)}3^{v_3(n)}5^{v_5(n)}
		\]
		und den zugehörigen rauen Kern durch
		\[
		m_{235}(n):=\frac{n}{s_{235}(n)}.
		\]
		Dann gilt
		\[
		n=s_{235}(n)\,m_{235}(n)
		\]
		sowie
		\[
		\gcd(m_{235}(n),30)=1.
		\]
	\end{definition}
	
	\begin{bemerkung}
		Der Faktor \(s_{235}(n)\) misst, wie stark \(n\) bereits durch die kleinen glatten
		Primträger \(2\), \(3\) und \(5\) aufgebaut wird. Der Faktor \(m_{235}(n)\) misst
		dagegen die neue Primstruktur, die nach Abzug dieser glatten Basis verbleibt.
	\end{bemerkung}
	
	\begin{satz}[Zerlegung in glatte Lage und rauhe Signatur]
		\label{satz:glatt_rau}
		Jede natürliche Zahl \(n\) besitzt bezüglich der Primträger \(2,3,5\) eine eindeutige
		Zerlegung
		\[
		\boxed{
			n=s_{235}(n)\,m_{235}(n).
		}
		\]
		Dabei gilt:
		
		\begin{enumerate}
			\item \(s_{235}(n)\) bestimmt die glatte Schalenlage von \(n\),
			\item \(m_{235}(n)\) bestimmt die rauhe Reststruktur,
			\item die feine Signatur des rauen Kerns wird durch
			\[
			m_{235}(n)\bmod 12 \in \{1,5,7,11\}
			\]
			und damit durch die Familien
			\[
			E,\ A,\ B,\ C
			\]
			beschrieben,
			\item die grobe Signatur des rauen Kerns wird durch die Zweischale
			\[
			EA\mid BC
			\]
			beschrieben.
		\end{enumerate}
	\end{satz}
	
	\begin{proof}
		Die Eindeutigkeit der Zerlegung folgt unmittelbar aus der eindeutigen Primfaktorzerlegung.
		Die Faktoren \(2\), \(3\) und \(5\) werden vollständig in \(s_{235}(n)\) gesammelt, sodass
		der Rest \(m_{235}(n)\) zu \(30\) teilerfremd ist. Daher können seine Primfaktoren modulo
		\(12\) nur in den Klassen
		\[
		1,\ 5,\ 7,\ 11
		\]
		liegen, also genau in \(E,A,B,C\). Die grobe EA/BC-Signatur ist die entsprechende
		Paarprojektion.
	\end{proof}
	
	\begin{definition}[Feine und grobe Signatur des rauen Kerns]
		Für \(m:=m_{235}(n)\) definieren wir die feine Signatur
		\[
		\sigma_4(m)\in\{E,A,B,C\}
		\]
		durch
		\[
		m\bmod 12=
		\begin{cases}
			1,& \sigma_4(m)=E,\\
			5,& \sigma_4(m)=A,\\
			7,& \sigma_4(m)=B,\\
			11,& \sigma_4(m)=C,
		\end{cases}
		\]
		sowie die grobe Signatur
		\[
		\chi_{EA/BC}(m)\in\{+1,-1\}
		\]
		durch
		\[
		\chi_{EA/BC}(m)=
		\begin{cases}
			+1,& m\bmod 12\in\{1,5\},\\[1mm]
			-1,& m\bmod 12\in\{7,11\}.
		\end{cases}
		\]
	\end{definition}
	
	\begin{korollar}[Natürliche kombinierte Schalenkoordinate]
		\label{kor:kombinierte_koordinate}
		Jede natürliche Zahl \(n\) besitzt damit eine natürliche kombinierte Schalenkoordinate der
		Form
		\[
		\boxed{
			\bigl(s_{235}(n),\,\sigma_4(m_{235}(n))\bigr)
		}
		\]
		bzw. in grober Form
		\[
		\boxed{
			\bigl(s_{235}(n),\,\chi_{EA/BC}(m_{235}(n))\bigr).
		}
		\]
	\end{korollar}
	
	\begin{bemerkung}
		Diese Koordinaten haben eine klare Interpretation:
		
		\begin{itemize}
			\item \(s_{235}(n)\) beschreibt die radiale Lage in der glatten Schale,
			\item \(\sigma_4(m_{235}(n))\) beschreibt die feine sektorielle Orientierung des rauen
			Kerns,
			\item \(\chi_{EA/BC}(m_{235}(n))\) beschreibt die grobe Paarorientierung des rauen Kerns.
		\end{itemize}
		
		Damit erscheint jede natürliche Zahl als Punkt in einer diskreten Zweiebenen-Geometrie:
		\[
		\text{glatte Lage} \times \text{raue Signatur}.
		\]
	\end{bemerkung}
	
	\begin{satz}[Schalenabschluss und raue Signatur]
		\label{satz:glatt_rau_abschluss}
		Für die Schale \(\{2,3,5\}\) gilt:
		\[
		m_{235}(n)=1
		\quad\Longleftrightarrow\quad
		n \text{ ist vollständig }(2,3,5)\text{-glatt.}
		\]
		Dagegen gilt:
		\[
		m_{235}(n)>1
		\quad\Longleftrightarrow\quad
		\text{in }n\text{ tritt neue Primstruktur jenseits der glatten Schale auf.}
		\]
		In diesem Fall beschreiben \(\sigma_4(m_{235}(n))\) und \(\chi_{EA/BC}(m_{235}(n))\)
		die innere Signatur dieses neuen Restes.
	\end{satz}
	
	\begin{proof}
		Dies ist der Spezialfall des allgemeinen Schalenabschluss-Kriteriums, angewandt auf die
		Primschale \(\{2,3,5\}\). Die Signaturaussagen folgen aus der modulo-\(12\)-Klassifikation
		des rauen Kerns.
	\end{proof}
	
	\begin{bemerkung}
		Damit ergibt sich eine natürliche Rollenverteilung:
		
		\[
		\boxed{
			\text{glatte Schale = Träger,}
			\qquad
			\text{rauer Kern = neue Struktur,}
			\qquad
			\text{EABC/EA-BC = Orientierung dieser neuen Struktur.}
		}
		\]
		
		Die glatten Zahlen geben die radiale Schalenordnung vor, die rauhen Signaturen liefern
		die lokale und globale sektorielle Differenzierung.
	\end{bemerkung}
	
	\begin{korollar}[Primzahlvierlinge im glatt-rauen Bild]
		\label{kor:vierling_glatt_rau}
		Für einen Primzahlvierling
		\[
		(p,p+2,p+6,p+8)
		\]
		mit \(p>5\) ist der glatte Anteil trivial, also
		\[
		s_{235}(p)=s_{235}(p+2)=s_{235}(p+6)=s_{235}(p+8)=1,
		\]
		sodass die gesamte Struktur bereits im rauen Kern liegt. Damit tragen Primzahlvierlinge
		im glatt-rauen Bild unmittelbar ihre volle feine EABC-Signatur bzw. ihre grobe EA/BC-
		Projektion.
	\end{korollar}
	
	\begin{bemerkung}
		Gerade deshalb ergänzen sich glatte Schalen und Primzahlvierlinge so gut: Die glatten
		Anteile organisieren die globale radiale Lage der Zahlen, während Primzahlvierlinge eine
		lokal maximal dichte Form reiner rauer Signatur darstellen.
	\end{bemerkung}
	
	\subsection{Numerischer Vierling-\(\to\)-Fünfling-Balance-Test}
	
	Um die Hypothese eines strukturellen Zusammenhangs zwischen dichten Primzahlvierlingen und
	dichten Primzahlfünflingen zu testen, wurde die Erweiterung
	\[
	Q_4(p):=(p,p+2,p+6,p+8)
	\quad\longrightarrow\quad
	Q_5(p):=(p,p+2,p+6,p+8,p+12)
	\]
	numerisch untersucht. Dabei wurden alle dichten Vierlinge bis \(10^6\) bestimmt und
	überprüft, welche davon durch das zusätzliche Glied \(p+12\) zu einem dichten
	Primzahlfünfling ergänzt werden.
	
	Für feste Pivot-Primzahlen \(\pi\) wurden die groben Balancegrößen
	\[
	\Delta_\pi,\qquad \Sigma_\pi
	\]
	sowie die logarithmische Balance
	\[
	B_\pi
	\qquad\text{und}\qquad
	\beta_\pi:=\frac{B_\pi}{\log 5}
	\]
	für den Vierling und den zugehörigen Fünfling verglichen.
	
	\begin{satz}[Numerischer Vierling-\(\to\)-Fünfling-Balance-Test]
		\label{satz:vierling_fuenfling_balance}
		Bis \(10^6\) wurden
		\[
		166
		\]
		dichte Primzahlvierlinge gefunden, von denen
		\[
		34
		\]
		zu dichten Primzahlfünflingen ergänzt werden konnten. Dies entspricht einer
		Ergänzungsrate von
		\[
		\frac{34}{166}\approx 20.48\%.
		\]
		
		Für die Pivot-Primzahlen
		\[
		\pi\in\{11,13,17,19,29,37\}
		\]
		zeigen die numerischen Daten:
		
		\begin{enumerate}
			\item Das fünfte Glied \(p+12\) des Fünflings liegt fast immer oberhalb des Pivots
			\(\pi\).
			
			\item Die diskrete Unter-/Ober-Balance verschiebt sich beim Übergang
			\[
			Q_4(p)\to Q_5(p)
			\]
			nahezu deterministisch um
			\[
			\Delta_\pi(Q_5)-\Delta_\pi(Q_4)\approx -1.
			\]
			
			\item Die gewichtete Zweischalen-Balance \(\Sigma_\pi\) reagiert auf die EA/BC-Signatur des
			neu hinzugefügten fünften Glieds.
			
			\item Die logarithmische Balance \(B_\pi\) erfährt einen starken negativen Sprung,
			während ein echter Vorzeichenwechsel von \(B_\pi\) nur selten auftritt.
		\end{enumerate}
	\end{satz}
	
	\begin{proof}[Numerische Verifikation]
		Die numerische Auswertung ergab für alle getesteten Pivot-Primzahlen dieselbe
		Ergänzungsrate
		\[
		20.48\%.
		\]
		
		Der Anteil der Fälle, in denen das neue fünfte Glied \(p+12\) oberhalb des Pivots liegt,
		beträgt:
		
		\[
		100.00\%\quad\text{für }\pi=11,13,
		\]
		\[
		97.06\%\quad\text{für }\pi=17,19,
		\]
		\[
		94.12\%\quad\text{für }\pi=29,37.
		\]
		
		Die mittlere Änderung der diskreten Balance \(\Delta_\pi\) beträgt numerisch:
		
		\[
		-1.000,\ -1.000,\ -0.971,\ -0.941,\ -0.882,\ -0.882,
		\]
		also durchgehend nahe bei \(-1\).
		
		Die mittlere Änderung der logarithmischen Balance \(B_\pi\) beträgt numerisch:
		
		\[
		-11.448843,\ -11.448843,\ -11.365514,\ -11.282184,\ -11.097743,\ -11.097743,
		\]
		und entsprechend
		\[
		\Delta\beta_\pi\approx -7.11,\ -7.11,\ -7.06,\ -7.01,\ -6.90,\ -6.90.
		\]
		
		Dagegen bleibt der Anteil echter Vorzeichenwechsel von \(B_\pi\) sehr klein:
		\[
		2.94\%,\ 0.00\%,\ 2.94\%,\ 0.00\%,\ 0.00\%,\ 0.00\%.
		\]
	\end{proof}
	
	\begin{bemerkung}
		Die numerischen Resultate zeigen, dass die Ergänzung eines dichten Vierlings zu einem
		dichten Fünfling typischerweise keine Spiegelung der Balance darstellt, sondern eine
		gerichtete Translation. Das neue fünfte Glied tritt fast immer oberhalb des Pivots auf
		und verschiebt dadurch die Unter-/Ober-Balance schichtweise nach oben.
	\end{bemerkung}
	
	\begin{bemerkung}
		Insbesondere wird damit die naive Formulierung
		\[
		\text{„die }5\text{ springt von unten nach oben“}
		\]
		in strenger Vorzeichenwechsel-Form nicht bestätigt. Bestätigt wird jedoch eine schwächere,
		aber strukturierte Aussage:
		\[
		Q_4\to Q_5
		\]
		wirkt numerisch als kontrollierter oberer Balance-Shift mit stabiler Änderung der
		diskreten Balance und mit EA/BC-kodierter Reaktion in \(\Sigma_\pi\).
	\end{bemerkung}
	
	\begin{korollar}[Grobes Zweischalenmuster der Erweiterung]
		\label{kor:vierling_fuenfling_grobmuster}
		Die groben EA/BC-Muster der getesteten Erweiterungen haben durchgehend die Form
		
		\[
		EA-BC-BC-EA \;\longrightarrow\; EA-BC-BC-EA-EA
		\]
		oder
		\[
		BC-EA-EA-BC \;\longrightarrow\; BC-EA-EA-BC-BC.
		\]
		
		Damit erscheint der dichte Primzahlfünfling numerisch als sektorkompatibler Randaufsatz
		eines dichten Primzahlvierlings.
	\end{korollar}
	
	\begin{proof}[Numerische Beobachtung]
		Die ausgegebenen Beispielmuster zeigen genau diese beiden Typen. Das zusätzliche fünfte
		Glied trägt jeweils dieselbe grobe EA/BC-Sektorart wie der entsprechende Rand des
		Vierlings.
	\end{proof}
	
	\begin{bemerkung}
		Die Ergebnisse verbinden die statische Quotientenschale
		\[
		EA\mid BC
		\]
		mit einer dynamischen Erweiterungsregel: Die Zweischale ist nicht nur als grobe Projektion
		der feinen Viererschale numerisch stabil, sondern bleibt auch beim Übergang von
		Vierlingen zu Fünflingen strukturell sichtbar.
	\end{bemerkung}
	
	\begin{bemerkung}
		Die nächste natürliche Fragestellung besteht darin, die echten Fünflingserweiterungen mit
		den nicht ergänzbaren Vierlingen zu vergleichen. Erst ein solcher Kontrasttest kann
		klären, ob die Ergänzbarkeit selbst durch eine charakteristische Balance- oder
		Zweischalenlage des Vierlings vorhergesagt wird.
	\end{bemerkung}
	
	\subsection{Musterspezifische Balance-Selektion bei der Vierling-\(\to\)-Fünfling-Erweiterung}
	
	Die bisherigen numerischen Tests deuten darauf hin, dass die Ergänzbarkeit eines dichten
	Primzahlvierlings zu einem dichten Primzahlfünfling nicht durch eine globale Balancegröße
	allein beschrieben wird. Vielmehr zeigt sich eine musterspezifische Kopplung zwischen der
	groben EA/BC-Struktur des Vierlings und seiner logarithmischen Balance.
	
	Zur Untersuchung wurden die dichten Primzahlvierlinge
	\[
	Q_4(p):=(p,p+2,p+6,p+8)
	\]
	bis \(10^6\) betrachtet und nach ihrem groben Zweischalenmuster in die beiden Typen
	\[
	EA-BC-BC-EA
	\qquad\text{und}\qquad
	BC-EA-EA-BC
	\]
	zerlegt. Innerhalb jedes Typs wurden die Vierlinge separat nach der Balancegröße
	\[
	B_\pi(Q_4)
	\]
	in Quantile eingeteilt und die Ergänzungsrate zum dichten Primzahlfünfling
	\[
	Q_5(p):=(p,p+2,p+6,p+8,p+12)
	\]
	bestimmt.
	
	\begin{satz}[Musterspezifische Balance-Selektion]
		\label{satz:musterspezifische_balance_selektion}
		Für die dichten Primzahlvierlinge bis \(10^6\) und den Pivot \(\pi=11\) gilt numerisch:
		
		\begin{enumerate}
			\item Im Typ
			\[
			EA-BC-BC-EA
			\]
			steigt die Ergänzungsrate mit wachsender Balance \(B_\pi(Q_4)\) deutlich an.
			
			\item Im Typ
			\[
			BC-EA-EA-BC
			\]
			ist dieser Zusammenhang deutlich schwächer und nicht monoton stabil.
			
			\item Die Ergänzbarkeit zu einem dichten Primzahlfünfling ist daher nicht durch
			\(B_\pi(Q_4)\) allein, sondern durch eine Kopplung der Form
			\[
			\boxed{
				\text{Mustertyp} \times \text{Balancehöhe}
			}
			\]
			bestimmt.
		\end{enumerate}
	\end{satz}
	
	\begin{proof}[Numerische Verifikation]
		Für den Typ
		\[
		BC-EA-EA-BC
		\]
		ergaben die fünf Balance-Quantile die Ergänzungsraten
		\[
		25.00\%,\ 12.50\%,\ 11.76\%,\ 12.50\%,\ 23.53\%.
		\]
		Dies zeigt zwar eine gewisse Variation, aber keinen klaren monotonen Anstieg mit wachsender
		Balance.
		
		Für den Typ
		\[
		EA-BC-BC-EA
		\]
		ergaben dieselben Quantile dagegen
		\[
		12.50\%,\ 23.53\%,\ 17.65\%,\ 29.41\%,\ 35.29\%.
		\]
		Hier ist eine deutliche Tendenz sichtbar: Vierlinge mit höherem \(B_\pi\)-Wert sind
		numerisch erheblich häufiger ergänzbar als Vierlinge mit niedriger Balance.
		
		Besonders aussagekräftig ist der Vergleich der Extremquantile:
		\[
		35.29\%-12.50\%=22.79\%.
		\]
		Die Differenz zwischen niedrigem und hohem Balancebereich ist damit im Typ
		\[
		EA-BC-BC-EA
		\]
		erheblich größer als im Typ
		\[
		BC-EA-EA-BC.
		\]
	\end{proof}
	
	\begin{bemerkung}
		Die globale Korrelation zwischen \(B_\pi(Q_4)\) und der Ergänzbarkeit war zuvor nur
		moderat. Der getrennte Quantiltest erklärt dies: Das Gesamtsignal mittelt über zwei
		unterschiedliche Kanäle, von denen nur einer stark balanceempfindlich ist.
	\end{bemerkung}
	
	\begin{korollar}[Balanceaktiver Ergänzungskanal]
		\label{kor:balanceaktiver_kanal}
		Der grobe Vierlingskanal
		\[
		EA-BC-BC-EA
		\]
		ist numerisch der eigentliche balanceaktive Ergänzungskanal. In diesem Kanal trägt ein
		hoher Wert von \(B_\pi(Q_4)\) deutlich zur späteren Ergänzbarkeit zum Fünfling bei.
	\end{korollar}
	
	\begin{bemerkung}
		Dagegen erscheint der Kanal
		\[
		BC-EA-EA-BC
		\]
		deutlich träger. Hier ist die Ergänzbarkeit wesentlich schwächer an die Balancehöhe
		gekoppelt. Dies spricht für zwei verschieden empfindliche sektorale Erweiterungskanäle.
	\end{bemerkung}
	
	\begin{satz}[Stabilität des 2D-Befunds über mehrere Pivots]
		\label{satz:2d_befund_mehrere_pivots}
		Die zweidimensionale Struktur
		\[
		\text{Mustertyp} \times \text{Balancehöhe}
		\]
		bleibt numerisch stabil für die getesteten Pivot-Primzahlen
		\[
		\pi\in\{11,13,17,19,29,37\}.
		\]
		Insbesondere bleibt die Differenz zwischen hohem und niedrigem Balancebereich im Typ
		\[
		EA-BC-BC-EA
		\]
		durchgehend deutlich größer als im Typ
		\[
		BC-EA-EA-BC.
		\]
	\end{satz}
	
	\begin{proof}[Numerische Verifikation]
		Im Mehrpivot-2D-Test ergab sich für alle getesteten Pivots:
		
		Für den Typ
		\[
		BC-EA-EA-BC:
		\qquad
		18.92\%-15.56\%=3.36\%.
		\]
		
		Für den Typ
		\[
		EA-BC-BC-EA:
		\qquad
		30.43\%-15.79\%=14.65\%.
		\]
		
		Diese Differenzen blieben über alle getesteten Pivots gleich. Damit ist die Asymmetrie der
		beiden Kanäle numerisch robust.
	\end{proof}
	
	\begin{bemerkung}
		Die Pivot-Primzahl \(\pi\) scheint damit weniger als feinabzustimmender Strukturparameter zu
		wirken, sondern vielmehr als Messfenster für eine bereits im Vierling vorhandene
		Balanceeigenschaft. Die wesentliche Selektion sitzt in der Vierlingsstruktur selbst.
	\end{bemerkung}
	
	\begin{korollar}[Vorläufige Arbeitsvermutung]
		\label{kor:arbeitsvermutung_balance_kanal}
		Die Ergänzbarkeit dichter Primzahlvierlinge zu dichten Primzahlfünflingen folgt
		numerisch keiner allgemeinen globalen Balance-Selektion, sondern einer
		\emph{musterspezifischen Balance-Selektion}. Der Kanal
		\[
		EA-BC-BC-EA
		\]
		ist der dominante balanceaktive Erweiterungskanal.
	\end{korollar}
	
	\begin{bemerkung}
		Die nächste natürliche Aufgabe besteht darin, diese kanalspezifische Balance-Selektion
		noch feiner aufzulösen, etwa durch weitere Quantilunterteilungen im aktiven Kanal oder
		durch Hinzunahme zusätzlicher Strukturgrößen für die wenigen verbleibenden
		Gegenbeispiele.
	\end{bemerkung}
	
	\subsection{Ein kanal-korrigierter Einzeltrigger}
	
	Die bisherigen numerischen Tests haben gezeigt, dass die Ergänzbarkeit eines dichten
	Primzahlvierlings zu einem dichten Primzahlfünfling weder durch die Balancegröße
	\(B_\pi(Q_4)\) allein noch durch den groben Mustertyp allein vollständig beschrieben wird.
	Dagegen deutet sich an, dass eine geeignete \emph{kanal-korrigierte} Einzelgröße beide
	Aspekte in einer einzigen Triggerfunktion zusammenführen kann.
	
	\begin{definition}[Kanal-korrigierter Einzeltrigger]
		\label{def:kanal_korrigierter_trigger}
		Für einen dichten Primzahlvierling
		\[
		Q_4(p)=(p,p+2,p+6,p+8)
		\]
		und einen festen Pivot \(\pi\) definieren wir den additiven Einzeltrigger
		\[
		\boxed{
			T_c(Q_4):=B_\pi(Q_4)+c\cdot \mathbf 1_{\{EA-BC-BC-EA\}}(Q_4),
		}
		\]
		wobei
		\[
		\mathbf 1_{\{EA-BC-BC-EA\}}(Q_4)=
		\begin{cases}
			1,& \text{falls }Q_4\text{ vom groben Typ }EA-BC-BC-EA\text{ ist},\\
			0,& \text{falls }Q_4\text{ vom groben Typ }BC-EA-EA-BC\text{ ist},
		\end{cases}
		\]
		und \(c\ge 0\) ein numerisch zu kalibrierender Kanalbonus ist.
	\end{definition}
	
	\begin{bemerkung}
		Der Trigger \(T_c\) ist eine einzelne reelle Größe. Er erhält die Balanceinformation
		\(B_\pi(Q_4)\) für alle Vierlinge und verleiht zugleich dem numerisch aktiveren Kanal
		\[
		EA-BC-BC-EA
		\]
		einen festen additiven Bonus.
	\end{bemerkung}
	
	\begin{satz}[Numerische Verbesserung durch kanal-korrigierten Trigger]
		\label{satz:kanaltrigger_numerisch}
		Für die dichten Primzahlvierlinge bis \(10^6\) und den Pivot \(\pi=11\) verbessert der
		kanal-korrigierte Trigger \(T_c\) die lineare und monotone Korrelation zur Ergänzbarkeit
		numerisch gegenüber dem reinen Balance-Trigger \(B_\pi(Q_4)\), sofern \(c\) im Bereich
		mittlerer positiver Werte gewählt wird.
	\end{satz}
	
	\begin{proof}[Numerische Verifikation]
		Für den reinen Balance-Trigger \(T_0(Q_4)=B_\pi(Q_4)\) ergaben sich die Korrelationen
		\[
		\mathrm{corr}_{\mathrm{Pearson}}(T_0,\text{Ergänzbarkeit})=0.200216,
		\]
		\[
		\mathrm{corr}_{\mathrm{Spearman}}(T_0,\text{Ergänzbarkeit})=0.117749.
		\]
		
		Für verschiedene Werte von \(c\) ergaben sich numerisch:
		
		\[
		c=2:
		\qquad
		0.206877,\ 0.128963,
		\]
		
		\[
		c=5:
		\qquad
		0.211361,\ 0.143915,
		\]
		
		\[
		c=10:
		\qquad
		0.207735,\ 0.146407,
		\]
		
		\[
		c=15:
		\qquad
		0.197285,\ 0.129897.
		\]
		
		Damit wird der Trigger im Bereich
		\[
		c\approx 5\text{ bis }10
		\]
		klar besser als \(B_\pi\) allein, während zu große Werte von \(c\) die Gesamtordnung wieder
		verschlechtern.
	\end{proof}
	
	\begin{korollar}[Bester numerischer Kandidatenbereich]
		\label{kor:bester_triggerbereich}
		Der bislang beste Kandidat für einen einzelnen numerischen Trigger der Form
		\[
		T_c(Q_4)=B_\pi(Q_4)+c\cdot \mathbf 1_{\{EA-BC-BC-EA\}}(Q_4)
		\]
		liegt im Bereich
		\[
		\boxed{
			c\in[5,10].
		}
		\]
	\end{korollar}
	
	\begin{bemerkung}
		Der Wert \(c=5\) liefert die beste lineare Korrelation, während \(c=10\) die beste
		Rangkorrelation erzeugt. Dies spricht dafür, dass ein moderater Kanalbonus den aktiven
		Vierlingskanal stärkt, ohne die Ordnung der Gesamtmenge zu stark zu verzerren.
	\end{bemerkung}
	
	\begin{bemerkung}
		Ein zuvor getesteter multiplikativer Trigger der Form
		\[
		\mathbf 1_{\{EA-BC-BC-EA\}}(Q_4)\cdot B_\pi(Q_4)
		\]
		erwies sich dagegen als ungeeignet. Der Grund liegt in der Entstehung einer großen
		Nullmasse für den zweiten Kanal, die die Rangordnung der Vierlinge zerstört und die
		Korrelation praktisch auf null reduziert. Der additive Trigger vermeidet dieses Problem.
	\end{bemerkung}
	
	\begin{satz}[Interpretation des Kanalbonus]
		\label{satz:interpretation_kanalbonus}
		Der additive Trigger \(T_c\) lässt sich so interpretieren, dass die Balanceinformation
		\(B_\pi(Q_4)\) prinzipiell für alle Vierlinge relevant bleibt, im Kanal
		\[
		EA-BC-BC-EA
		\]
		jedoch mit systematisch höherer Wirksamkeit in Bezug auf die Fünflingsergänzbarkeit
		auftritt.
	\end{satz}
	
	\begin{proof}[Heuristische Interpretation]
		Die getrennten Quantil- und 2D-Tests zeigten, dass derselbe Anstieg der Balancehöhe in den
		beiden groben Vierlingsmustern sehr unterschiedlich wirkt. Im Kanal
		\[
		EA-BC-BC-EA
		\]
		ist der Zusammenhang zwischen Balancehöhe und Ergänzbarkeit deutlich stärker als im Kanal
		\[
		BC-EA-EA-BC.
		\]
		Ein additiver Bonus ist daher die einfachste Einzelgröße, welche diese Asymmetrie
		modelliert, ohne die Gesamtordnung der Balanceachse zu zerstören.
	\end{proof}
	
	\begin{korollar}[Vorläufige Einzeltrigger-Vermutung]
		\label{kor:einzeltrigger_vermutung}
		Ein numerisch plausibler Einzeltrigger für die Ergänzbarkeit dichter Primzahlvierlinge zu
		dichten Primzahlfünflingen ist gegeben durch
		\[
		\boxed{
			T^*(Q_4)=B_\pi(Q_4)+c\cdot \mathbf 1_{\{EA-BC-BC-EA\}}(Q_4),
			\qquad c\approx 5\text{ bis }10.
		}
		\]
		Große Werte von \(T^*(Q_4)\) korrelieren numerisch stärker mit der Ergänzbarkeit als die
		reine Balancegröße \(B_\pi(Q_4)\) allein.
	\end{korollar}
	
	\begin{bemerkung}
		Diese Vermutung ist derzeit rein numerisch motiviert. Der nächste natürliche Schritt
		besteht darin, den optimalen Bereich für \(c\) feiner aufzulösen und zu untersuchen, ob
		sich der Parameter aus einer tieferen strukturellen Eigenschaft der groben EA/BC-Schale
		ableiten lässt.
	\end{bemerkung}
	
	\subsection{Feinkalibrierung des Kanalbonus und strukturelle Deutung von \(c^*=6\)}
	
	Nachdem sich gezeigt hatte, dass ein additiver kanal-korrigierter Trigger die
	Ergänzbarkeit dichter Primzahlvierlinge besser beschreibt als die reine Balancegröße
	\(B_\pi(Q_4)\), wurde der Kanalbonus \(c\) feiner kalibriert. Untersucht wurde die
	Einzelgröße
	\[
	T_c(Q_4):=B_\pi(Q_4)+c\cdot \mathbf 1_{\{EA-BC-BC-EA\}}(Q_4).
	\]
	
	Dabei bezeichnet
	\[
	\mathbf 1_{\{EA-BC-BC-EA\}}(Q_4)
	\]
	den Indikator des numerisch aktiven Vierlingskanals
	\[
	EA-BC-BC-EA.
	\]
	
	\begin{satz}[Feinkalibrierung des Kanalbonus]
		\label{satz:feinkalibrierung_kanalbonus}
		Für die dichten Primzahlvierlinge bis \(10^6\) und den Pivot \(\pi=11\) besitzt der
		additive Trigger
		\[
		T_c(Q_4)=B_\pi(Q_4)+c\cdot \mathbf 1_{\{EA-BC-BC-EA\}}(Q_4)
		\]
		ein stabiles numerisches Optimum im Bereich
		\[
		c\approx 6.
		\]
		Genauer liefert \(c=6\) die beste lineare Korrelation zur Ergänzbarkeit, während die
		Rangkorrelation ihr Optimum erst bei etwas größeren Werten, insbesondere nahe \(c=9\),
		erreicht.
	\end{satz}
	
	\begin{proof}[Numerische Verifikation]
		Im Feinscan für
		\[
		c\in\{5.0,5.5,6.0,6.5,7.0,7.5,8.0,8.5,9.0\}
		\]
		ergaben sich die folgenden Werte:
		
		\[
		c=5.0:
		\qquad
		\mathrm{corr}_{\mathrm{Pearson}}=0.211361,
		\qquad
		\mathrm{corr}_{\mathrm{Spearman}}=0.143915,
		\]
		
		\[
		c=5.5:
		\qquad
		0.211535,
		\qquad
		0.142046,
		\]
		
		\[
		c=6.0:
		\qquad
		0.211565,
		\qquad
		0.141423,
		\]
		
		\[
		c=6.5:
		\qquad
		0.211462,
		\qquad
		0.142046,
		\]
		
		\[
		c=7.0:
		\qquad
		0.211233,
		\qquad
		0.143915,
		\]
		
		\[
		c=7.5:
		\qquad
		0.210888,
		\qquad
		0.144849,
		\]
		
		\[
		c=8.0:
		\qquad
		0.210437,
		\qquad
		0.145472,
		\]
		
		\[
		c=8.5:
		\qquad
		0.209887,
		\qquad
		0.147965,
		\]
		
		\[
		c=9.0:
		\qquad
		0.209248,
		\qquad
		0.149522.
		\]
		
		Daraus folgt:
		
		\begin{enumerate}
			\item Das Maximum der linearen Korrelation liegt bei
			\[
			\boxed{c=6.0.}
			\]
			
			\item Das Maximum der Rangkorrelation liegt bei
			\[
			\boxed{c=9.0.}
			\]
			
			\item Im Bereich
			\[
			5.5\le c\le 6.5
			\]
			bleibt der Pearson-Wert praktisch konstant auf seinem Maximalniveau. Dies zeigt, dass
			das Optimum nicht nur punktuell, sondern als stabiles lokales Fenster um \(6\) herum
			auftritt.
		\end{enumerate}
	\end{proof}
	
	\begin{korollar}[Vorläufiger Hauptwert des Kanalbonus]
		\label{kor:hauptwert_kanalbonus}
		Als numerisch und strukturell bevorzugter Hauptwert des Kanalbonus kann vorläufig
		\[
		\boxed{c^*=6}
		\]
		gewählt werden.
	\end{korollar}
	
	\begin{bemerkung}
		Die Wahl \(c^*=6\) wird nicht allein durch die numerische Maximierung gestützt, sondern
		auch durch ihre strukturelle Plausibilität im Vierlingskontext.
	\end{bemerkung}
	
	\begin{satz}[Strukturelle Plausibilität der Zahl \(6\)]
		\label{satz:strukturelle_plausibilitaet_6}
		Die Zahl \(6\) ist im Kontext dichter Primzahlvierlinge in mehrfacher Hinsicht natürlich:
		
		\begin{enumerate}
			\item Sie ist die kleinste volle glatte Grundperiode
			\[
			6=2\cdot 3,
			\]
			relativ zu der Primzahlen \(>3\) stabil als Restklassen sichtbar werden.
			
			\item Im Vierling
			\[
			Q_4(p)=(p,p+2,p+6,p+8)
			\]
			trennt sie die beiden inneren Primzahlpaare
			\[
			(p,p+2)
			\qquad\text{und}\qquad
			(p+6,p+8).
			\]
			
			\item Sie kann daher als natürliche Transport-, Synchronisations- oder
			Kanalverschiebungslänge des aktiven Erweiterungskanals interpretiert werden.
		\end{enumerate}
	\end{satz}
	
	\begin{proof}[Heuristische Interpretation]
		Die dichten Primzahlvierlinge bestehen aus zwei eng gekoppelten Zweierpaaren:
		\[
		(p,p+2)
		\qquad\text{und}\qquad
		(p+6,p+8).
		\]
		Diese beiden Paare sind gegeneinander genau um
		\[
		6
		\]
		verschoben. Zugleich ist \(6\) die kleinste nichttriviale \(2,3\)-glatte Grundperiode, in
		der Primzahlmuster jenseits von \(2\) und \(3\) stabil auftreten.
		
		Wenn der Kanal
		\[
		EA-BC-BC-EA
		\]
		tatsächlich der balanceaktive Erweiterungskanal ist, dann liegt es nahe, dass sein
		optimierender Bonus dieselbe Größenordnung besitzt wie diese minimale innere
		Vierlingsverschiebung. Der numerische Befund
		\[
		c^*\approx 6
		\]
		passt auffallend gut zu dieser Struktur.
	\end{proof}
	
	\begin{bemerkung}
		Diese Deutung ist derzeit heuristisch. Sie zeigt jedoch, dass die numerische Kalibrierung
		des Triggerparameters nicht losgelöst von der inneren Geometrie des Vierlings gelesen
		werden sollte.
	\end{bemerkung}
	
	\begin{korollar}[Vorläufige Triggerform]
		\label{kor:vorlaeufige_triggerform}
		Der derzeit beste einfache numerische Einzeltrigger für die Ergänzbarkeit dichter
		Primzahlvierlinge zu dichten Primzahlfünflingen ist gegeben durch
		\[
		\boxed{
			T^*(Q_4)=B_\pi(Q_4)+6\cdot \mathbf 1_{\{EA-BC-BC-EA\}}(Q_4).
		}
		\]
	\end{korollar}
	
	\begin{bemerkung}
		Eine alternative, stärker rangorientierte Variante ist
		\[
		T^\sharp(Q_4)=B_\pi(Q_4)+9\cdot \mathbf 1_{\{EA-BC-BC-EA\}}(Q_4),
		\]
		die numerisch eine etwas höhere Spearman-Korrelation erzeugt. Für die theoretische
		Interpretation erscheint jedoch \(c^*=6\) aufgrund seiner strukturellen Natürlichkeit als
		der überzeugendere Hauptwert.
	\end{bemerkung}
	
	\begin{bemerkung}
		Der nächste natürliche Schritt besteht darin, die Stabilität des Optimums \(c^*=6\) über
		weitere Datenbereiche und zusätzliche Pivot-Primzahlen hinweg zu prüfen und zu
		untersuchen, ob der Bonus \(6\) aus einer tieferen EA/BC-Schalenstruktur formal
		abgeleitet werden kann.
	\end{bemerkung}
	
	\subsection{Korrelation des kanal-korrigierten Triggers mit der Distanz zum nächsten Primzahlzwilling}
	
	Nachdem sich gezeigt hatte, dass der kanal-korrigierte Trigger
	\[
	T^*(Q_4)=B_\pi(Q_4)+6\cdot \mathbf 1_{\{EA-BC-BC-EA\}}(Q_4)
	\]
	die Ergänzbarkeit dichter Primzahlvierlinge zu dichten Primzahlfünflingen besser beschreibt
	als die reine Balancegröße \(B_\pi(Q_4)\), stellt sich die Frage, ob dieselbe Größe auch mit
	der weiteren lokalen Primzahldichte zusammenhängt. Eine natürliche Zielgröße hierfür ist die
	Distanz vom Vierling zum nächsten Primzahlzwilling.
	
	\begin{definition}[Zwillingsdistanz]
		\label{def:zwillingsdistanz}
		Sei
		\[
		Q_4(p)=(p,p+2,p+6,p+8)
		\]
		ein dichter Primzahlvierling. Der \emph{nächste Primzahlzwilling rechts von \(Q_4\)} sei
		\[
		Z_{\mathrm{next}}(Q_4)=(q,q+2)
		\qquad\text{mit}\qquad q>p+8,
		\]
		wobei \(q\) minimal gewählt wird. Dann definieren wir die \emph{Zwillingsdistanz}
		\[
		\mathrm{Gap}(Q_4):=q-(p+8).
		\]
		Zusätzlich betrachten wir die logarithmische Distanz
		\[
		\log \mathrm{Gap}(Q_4).
		\]
	\end{definition}
	
	\begin{satz}[Der Trigger \(T^*\) als Distanzsignal]
		\label{satz:trigger_zwillingsdistanz}
		Für die dichten Primzahlvierlinge bis \(2\cdot 10^6\) korrelieren sowohl die Balancegröße
		\(B_\pi(Q_4)\) als auch der kanal-korrigierte Trigger \(T^*(Q_4)\) negativ mit der Distanz
		zum nächsten Primzahlzwilling. Der Trigger \(T^*(Q_4)\) trägt dabei das stärkere Signal.
	\end{satz}
	
	\begin{proof}[Numerische Verifikation]
		Für den Pivot \(\pi=11\) ergaben sich numerisch:
		
		\[
		\mathrm{corr}(B_\pi,\mathrm{Gap})=-0.186642,
		\qquad
		\mathrm{corr}(T^*,\mathrm{Gap})=-0.205656.
		\]
		
		Noch deutlicher wird der Zusammenhang auf logarithmischer Skala:
		
		\[
		\mathrm{corr}(B_\pi,\log \mathrm{Gap})=-0.267469,
		\qquad
		\mathrm{corr}(T^*,\log \mathrm{Gap})=-0.290561.
		\]
		
		Auch in der Rangordnung zeigt sich derselbe Effekt:
		
		\[
		\rho(B_\pi,\mathrm{Gap})=-0.113540,
		\qquad
		\rho(T^*,\mathrm{Gap})=-0.183009.
		\]
		
		Damit gilt numerisch:
		\[
		\boxed{
			\text{Je größer }T^*(Q_4),\text{ desto kleiner ist typischerweise die Distanz zum nächsten Primzahlzwilling.}
		}
		\]
	\end{proof}
	
	\begin{bemerkung}
		Die logarithmische Distanz \(\log \mathrm{Gap}\) reagiert stärker als die rohe Distanz
		\(\mathrm{Gap}\). Dies zeigt, dass der Trigger eher die Größenordnung der Zwillingsdistanz
		als einen festen additiven Abstand beschreibt.
	\end{bemerkung}
	
	\begin{korollar}[Der Trigger \(T^*\) ist ein Distanztrigger]
		\label{kor:distanztrigger}
		Der kanal-korrigierte Trigger \(T^*(Q_4)\) sollte nicht primär als Prädiktor für den Typ des
		nächsten Primzahlzwillings verstanden werden, sondern als numerisches Signal für die
		\emph{Nähe} des nächsten Zwillings.
	\end{korollar}
	
	\begin{satz}[Kanalasymmetrie der Zwillingsdistanz]
		\label{satz:kanalasymmetrie_zwillingsdistanz}
		Der grobe Vierlingskanal
		\[
		EA-BC-BC-EA
		\]
		liegt im Mittel näher am nächsten Primzahlzwilling als der Kanal
		\[
		BC-EA-EA-BC.
		\]
	\end{satz}
	
	\begin{proof}[Numerische Verifikation]
		Die Aufspaltung nach den beiden groben Vierlingsmustern ergab:
		
		Für
		\[
		BC-EA-EA-BC:
		\qquad
		\langle \mathrm{Gap}\rangle = 154.2517,
		\qquad
		\mathrm{Median}(\mathrm{Gap}) = 130.
		\]
		
		Für
		\[
		EA-BC-BC-EA:
		\qquad
		\langle \mathrm{Gap}\rangle = 130.5526,
		\qquad
		\mathrm{Median}(\mathrm{Gap}) = 88.
		\]
		
		Damit ist der numerisch aktive Vierlingskanal nicht nur balanceaktiver und
		fünflingsfreundlicher, sondern auch zwillingsnäher.
	\end{proof}
	
	\begin{bemerkung}
		Diese Asymmetrie ist insbesondere deshalb bemerkenswert, weil der Typ des nächsten
		Primzahlzwillings selbst kaum mit \(B_\pi\) oder \(T^*\) korreliert. Der Vierling beeinflusst
		also eher die \emph{Distanz} bis zur nächsten Zwillingslage als deren feinen Typ
		\((AB\) oder \(CE)\).
	\end{bemerkung}
	
	\begin{satz}[Quantilstruktur der Zwillingsdistanz]
		\label{satz:quantile_zwillingsdistanz}
		Sortiert man die Vierlinge nach dem kanal-korrigierten Trigger \(T^*(Q_4)\), so nimmt die
		mittlere und mediane Zwillingsdistanz über die Quantile hinweg deutlich ab.
	\end{satz}
	
	\begin{proof}[Numerische Verifikation]
		Für fünf Quantile nach \(T^*(Q_4)\) ergaben sich:
		
		\[
		Q_1:
		\qquad
		\langle \mathrm{Gap}\rangle = 169.4576,
		\qquad
		\mathrm{Median} = 148,
		\]
		
		\[
		Q_2:
		\qquad
		160.5085,
		\qquad
		118,
		\]
		
		\[
		Q_3:
		\qquad
		139.5593,
		\qquad
		112,
		\]
		
		\[
		Q_4:
		\qquad
		125.2203,
		\qquad
		88,
		\]
		
		\[
		Q_5:
		\qquad
		115.4576,
		\qquad
		82.
		\]
		
		Dies zeigt einen klaren abfallenden Trend. Insbesondere liegen Vierlinge mit hohem
		Triggerwert deutlich näher am nächsten Primzahlzwilling als Vierlinge mit niedrigem
		Triggerwert.
	\end{proof}
	
	\begin{bemerkung}
		Die entsprechende Quantilstruktur für \(B_\pi(Q_4)\) allein ist ebenfalls sichtbar, aber
		schwächer. Der Kanalbonus in \(T^*(Q_4)\) verbessert somit nicht nur die lokale
		Fünflingsergänzbarkeit, sondern auch die Erklärung der mittleren Zwillingsdistanz.
	\end{bemerkung}
	
	\begin{korollar}[Hierarchie der Wirkung von \(T^*\)]
		\label{kor:hierarchie_triggerwirkung}
		Der Trigger \(T^*(Q_4)\) wirkt numerisch auf mindestens zwei Ebenen:
		
		\begin{enumerate}
			\item lokal auf die Ergänzbarkeit
			\[
			Q_4\to Q_5,
			\]
			\item mittelfristig auf die Distanz
			\[
			Q_4\to Z_{\mathrm{next}}.
			\]
		\end{enumerate}
		
		Dagegen wirkt \(T^*(Q_4)\) kaum auf den feinen Typ des nächsten Zwillings.
	\end{korollar}
	
	\begin{bemerkung}
		Damit erscheint \(T^*(Q_4)\) weniger als Typtrigger, sondern vielmehr als eine Art
		\emph{Dichte- oder Distanztrigger}: Hohe Werte signalisieren eine Umgebung, in der der
		Zahlenraum statistisch schneller wieder in eine Zwillingslage zurückfindet.
	\end{bemerkung}
	
	\begin{bemerkung}
		Der nächste natürliche Schritt bestünde darin, dieselbe Frage nicht nur für den nächsten
		Primzahlzwilling, sondern auch für den nächsten dichten Primzahlvierling zu stellen. So
		ließe sich prüfen, ob der Trigger \(T^*(Q_4)\) allgemein mit lokaler Primzahlclusterung
		zusammenhängt.
	\end{bemerkung}
	
\end{document}